% -----------------------------------------------------------
\section{1840}
% -----------------------------------------------------------

	% ----------------------
	\subsection{Joseph Smith}
	% ----------------------
		\label{EMS-1840-JS}

		% ----------------------
		\subsubsection{Introduction to 1840}
		% ----------------------
		
			sdf
			
		% -----
		\subsubsection{Blessing to James Ivins, 7 January 1840}
		% -----
			\label{EMS-1840-JS-7-JAN-1840}
			% \label{EMS-1840-JS-DAY-3LETTERMONTH-YEAR}
			
			\begin{description}
				\item[Print Sources]
			\end{description}

			\begin{tabular}{p{0.5in}p{1in}p{0.5in}p{1in}}
				JSP	 	& ---		& JSR 		& --- \\
				DC	 	& ---		& HC 		& --- \\
				HDDC 	& ---		& TCHC 		& --- \\
			\end{tabular}
			% HDDC = woodford's DC diss; JSR = Marquardt's JS Revelations; TCHC = Vogel HC
			
			\begin{description}
				\item[Online Access] \href{https://www.josephsmithpapers.org/paper-summary/blessing-to-james-ivins-7-january-1840/1}{Here}
				% \item[Analysis] \hyperref[XX]{Here}
			\end{description}
			
			\begin{description}
				\item[Background]
			\end{description}
			
				% It might also be useful to give a two sentence context of the period: e.g., "This speech was given in Nauvoo to a small congregation one month prior to the destruction of the Nauvoo Expositor. These and other tensions may have been on the prophet's mind when he gave the speech."
			
			\begin{description}
				\item[Original Source]
			\end{description}
			
				% brief description of original source material, with reference to JSP or somewhere else for more info
				% Background material: it might be helpful to have a note that says whether a given document was presented as a speech, was written in a journal, etc.
				% Also a comment here about how reliable the source might be, or what shape it is in, if there are large omissions or parts that are hard to understand, and so on.
			
			\begin{description}
				\item[Text]
			\end{description}
			
				\begin{tightcenter}
						\hypertarget{EMS-1840-JS-7-JAN-1840-a}%
					Patria[r]chal
						\marginnote{a}%
					Blessing of James Ivins
				\end{tightcenter}

				Dear Brother, I lay my hands upon thy head in the name of Jesus Christ of Nazareth. and bless thee with a Patriarchal Blessing: in virtue of the Authority of the Holy Priesthood invested in me; and thou shall be blessed. Thy children also shall be blessed in the earth, and they shall min[i]ster to the children of men with power and Great glory, and thine eyes shall see it, thou shall be blessed with long life, and live to a good old age; and <thou> shall see thy childrens children unto the (\uline{second or third}) generation for thou hast a blessed family and they shall be blessed in the earth, unto the latest generation
					\hypertarget{EMS-1840-JS-7-JAN-1840-b}%
				and
					\marginnote{b}%
				thou shalt hear the voice of the Son of Man saying unto thee, Son they sins are forgiven thee And Angles Shalt minister unto thee, because of the honesty of thy heart, before the Lord; and it shall come to pass, that in after time the holy Preisthood shall be sealed upon thine head; and thou shall be blessed in thy basket and in thy store, and in thy soul, and in thy body, in thy goings out and in thy comings in. and the Lord will bless thee \uline{in the Land} and thou shall be found at the right hand of God. the Ancient of days shall stretch forth his hand to receive thee and thou shall minister before him, I now bless thee with \sout{the} \sout{blessing} the blessing of Abraham \& Isaac and of Jacob------

				and now I seal thee up \sout{into} <unto> eternal life; and pronounce these blessing upon thy head in the Name of Jesus Christ of Nazareth, even so Amen------

				Jan 7\supersu{th} 1840

				By Joseph Smith Jr, First President of the Church of Jesus Christ, of Latter Day Saints------
				
		% -----
		\subsubsection{Discourse, 13 January 1840}
		% -----
			\label{EMS-1840-JS-13-JAN-1840}
			% \label{EMS-1840-JS-DAY-3LETTERMONTH-YEAR}
			
			\begin{description}
				\item[Print Sources]
			\end{description}

			\begin{tabular}{p{0.5in}p{1in}p{0.5in}p{1in}}
				JSP	 	& ---		& JSR 		& --- \\
				DC	 	& ---		& HC 		& --- \\
				HDDC 	& ---		& TCHC 		& --- \\
			\end{tabular}
			% HDDC = woodford's DC diss; JSR = Marquardt's JS Revelations; TCHC = Vogel HC
			
			\begin{description}
				\item[Online Access] \href{https://www.josephsmithpapers.org/paper-summary/minutes-and-discourse-13-january-1840/1}{Here}
				% \item[Analysis] \hyperref[XX]{Here}
			\end{description}
			
			\begin{description}
				\item[Background]
			\end{description}
			
				% It might also be useful to give a two sentence context of the period: e.g., "This speech was given in Nauvoo to a small congregation one month prior to the destruction of the Nauvoo Expositor. These and other tensions may have been on the prophet's mind when he gave the speech."
			
			\begin{description}
				\item[Original Source]
			\end{description}
			
				% brief description of original source material, with reference to JSP or somewhere else for more info
				% Background material: it might be helpful to have a note that says whether a given document was presented as a speech, was written in a journal, etc.
				% Also a comment here about how reliable the source might be, or what shape it is in, if there are large omissions or parts that are hard to understand, and so on.
			
			\begin{description}
				\item[Text]
			\end{description}
			
				Brother Joseph Smith Jr dilated at some length on the offices of the Priesthood and on the duties of Elders Bishops. Priests. \&c and directed it should be entered on the minutes as the injunction of the Presidency that travelling Elders should be especially cautious of incroaching on the ground of stationed \& presiding Elders and rather direct their efforts to breaking up and occupying new ground and that the Churc[h]es generally refuse to be burthened with the support of unprofitable and dilatory labourers. It was unanimously resolved that this be received as the will and wish of the Conference. 
				
		% -----
		\subsubsection{Letter to Editor, 22 January 1840}
		% -----
			\label{EMS-1840-JS-22-JAN-1840}
			% \label{EMS-1840-JS-DAY-3LETTERMONTH-YEAR}
			
			\begin{description}
				\item[Print Sources]
			\end{description}

			\begin{tabular}{p{0.5in}p{1in}p{0.5in}p{1in}}
				JSP	 	& ---		& JSR 		& --- \\
				DC	 	& ---		& HC 		& --- \\
				HDDC 	& ---		& TCHC 		& --- \\
			\end{tabular}
			% HDDC = woodford's DC diss; JSR = Marquardt's JS Revelations; TCHC = Vogel HC
			
			\begin{description}
				\item[Online Access] \href{https://www.josephsmithpapers.org/paper-summary/letter-to-editor-22-january-1840/1}{Here}
				% \item[Analysis] \hyperref[XX]{Here}
			\end{description}
			
			\begin{description}
				\item[Background]
			\end{description}
			
				% It might also be useful to give a two sentence context of the period: e.g., "This speech was given in Nauvoo to a small congregation one month prior to the destruction of the Nauvoo Expositor. These and other tensions may have been on the prophet's mind when he gave the speech."
			
			\begin{description}
				\item[Original Source]
			\end{description}
			
				% brief description of original source material, with reference to JSP or somewhere else for more info
				% Background material: it might be helpful to have a note that says whether a given document was presented as a speech, was written in a journal, etc.
				% Also a comment here about how reliable the source might be, or what shape it is in, if there are large omissions or parts that are hard to understand, and so on.
			
			\begin{description}
				\item[Text]
			\end{description}
			
				\begin{tightright}
						\hypertarget{EMS-1840-JS-22-JAN-1840-a}%
					Brandywine
						\marginnote{a}%
					Chester Co Pa Jan. 22\supers{d} 1840
				\end{tightright}

				\sout{Faith of} the Latter Day Saints \sout{on Governments} \sout{and Laws in general}

				Mr Editor Sir

				For as much as many false rumors are a broad in the world concerning my self and the faith which I profess and that my belief with regard to Earthly governments and laws in general may not be miss interpreted nor miss understood I have thought proper to present for your consideration and for the consideration of the public (if you will do me the favour.) through your valuable and interesting paper my opinions concerning the same

					\hypertarget{EMS-1840-JS-22-JAN-1840-b}%
				First
					\marginnote{b}%
				I believe that goverments were instituted of God for the benefit of man and that he holds men accountable for \sout{their} their acts in relation to them Either in making laws or administering them for the good and Safety of Society

				<Secondly> I believe that no goverment can exist in peace except such laws are framed and held inviolate as will secure to each individual the free exercise of conscince the right and controll of property and the protection of life

				<Third> <I> \sout{We} believe that all governments necessarily require civel officers and magistrates to inforce the Laws of the Same and that such as will administer the Law \sout{<of the Same>} in equity and justice should be sought for and upheld by the voice of the people (if a republick) or the will of the Sovreign

					\hypertarget{EMS-1840-JS-22-JAN-1840-c}%
				<Fourth>
					\marginnote{c}%
				I believe that a religion is \sout{instiuted} instituted of God and that men are ameniable to him <and to him \sout{onley}> onley for the exercise of it unless their religious opinion prompts them to infringe upon the rights and <libertey> \sout{privalegs} \sout{<privaleges>} of others But I do not believe that human Law has a right to interfear in prescribing rules of worship to bind the conciences of men nor dictate forms for public or private devotion That the civel magistrate should restrane crime but never controll conscience should punish guilt but never supress the freedom of the <soul>

					\hypertarget{EMS-1840-JS-22-JAN-1840-d}%
				<5>
					\marginnote{d}%
				I believe that all men are bound to sustane and uphold the respective Governments in which they reside while protected in their inherent and in alienable rights by the Laws of such Governments and that Sedition and rebellion are unbecoming every Citizen thus prote[c]ted and should be punished accordingly and that all Governments have a right to enact such laws as\sout{s} in their own judgements are best calculated to secure the public interest at the same time <however> holding sacred the freedom of concience

					\hypertarget{EMS-1840-JS-22-JAN-1840-e}%
				<6>
					\marginnote{e}%
				I believe that every man should be honoured in his station Rulers and magistrates as such being plalaced [placed] for the protection of the inocent and the punishment of the guilty and that to the Laws all men owe respect and deference as with out them peace and harmony would be supplanted by anarchy and \sout{confusion} terror human Laws being instituted for the express purposs of regulating our interests as individuals and Nations between--- man and man and divine laws given of heaven prescribing rules on spiritual concerns for--- faith and worship both to be answered by man to his maker

					\hypertarget{EMS-1840-JS-22-JAN-1840-f}%
				<7>
					\marginnote{f}%
				I believe that rulers states and \sout{and} governments have a right and are bound to enact Laws for the protection of all Citizen in the free exercise of their religious belief But I do not believe that they have a right in justice to deprive Citizens of this privalege or proscribe them in their opinions so long as a regard and reverence \sout{are} \uline{is} Shown to the Laws and such religious opinions do not justify Sedition nor conspiracy

					\hypertarget{EMS-1840-JS-22-JAN-1840-g}%
				<8>
					\marginnote{g}%
				I\sout{do not} believe that the commission of crime should be punished according to the nature of the offence that murder treason Robbery theft and the breach of the general peace in all respects should be punished according to their criminalty and their tendancy to evil among men by the Laws of that Government in which the offence is committed and for the public peace and tranquility all men should step forward and use their ability in bringing \sout{the} offenders aggainst good laws to \sout{justice} punishment

				<9> I do not believe it just to mingle religious influence with civel government whereby one religious Society is fostered and an other proscribed in its spiritual privaleges and the individual rights of its members as citizens denied

					\hypertarget{EMS-1840-JS-22-JAN-1840-h}%
				<10>
					\marginnote{h}%
				I believe that all religious \sout{Society} Societies have a right to deal with \sout{its} <their> members for disorderly conduct according to the rules and regulations of such societies provided that such dealing be for fellowship and good Standing but \sout{we} <I> doo not believe that any religious Society has authority to try men \sout{for} <on> the right of property or life to take \sout{away} <from them> this world's goods or put them in Jepardy either life or limb neither to inflict any fisical [physical] punishment upon them They can onley excommunicate them from their society and \sout{with their} with draw \sout{their} <from their> fellowship

					\hypertarget{EMS-1840-JS-22-JAN-1840-i}%
				<11>
					\marginnote{i}%
				I believe that men should appeal to the Civel law for redress of all wrongs and grieveences where personal abuse is inflicted or the right of propperty or character infringed where such laws exist as will protect the Same but \sout{we} <I> believe that all men are Justified in defending thems Selves their friends and property and the government from <the> unlawful assaults and encroachments of all persons in times of exigincies where immediate appeal cannot be--- made to the Laws and relief afforded

					\hypertarget{EMS-1840-JS-22-JAN-1840-j}%
				<12>
					\marginnote{j}%
				I believe it just to preach the Gospel to the Nations of the Earth and warn the Righteous to Save themselves from the corruptions of the world But I do not believe it right to interfear with bond Servants neither preach the gospel to nor baptize them contrary to the will and wish of their masters nor to meddle with or influence them in the least \sout{contrary to the wish} to cause them to be dissatisfied with their Sittuations in this life theirby jeopardiseing the lives of men Such interfearence \sout{we} I believe to be unlawful and unjust and dangerous to the peace of every government allowing human beings to be held in Servitude

					\hypertarget{EMS-1840-JS-22-JAN-1840-k}%
				<13>
					\marginnote{k}%
				It has been reported by some vicious or de[s]igning characters that the church of Latter Day Saints believe in having their pro[p]erty in common and also the leaders of sade church controlls Said propperty--- This is a base fabrication without the least the least shadow or collering of any thing to make it out of but on the contrary no persons feelings can be more repugnant to such a principle than mine6 Every person in this Church has a right to controll his own propperty and is not required to do any thing except by his own free voluntary act That he may impart to the poor according to the requirement of the gospel. ``Give to him that asketh thee and from him that would borrow of thee turn <not> thou \sout{not} away'' \sout{I} Math 5 chap 42. v.

					\hypertarget{EMS-1840-JS-22-JAN-1840-l}%
				I
					\marginnote{l}%
				believe in liveing a virtuous upright and holy life before God and feel it my duty <to> perswad [persuade] all men in my power to do the same: That they may ciase [cease] to do evil and learn to do well and brake off \sout{from} their Sins by Rigteousness

				I close this by subscribing my self your\sout{s} mo[st] obedient Servent

				\begin{tightright}
					\textbf{Joseph Smith Jr}---
				\end{tightright}
			
		% -----
		\subsubsection{Letter to Emma Smith, 20--25 January 1840}
		% -----
			\label{EMS-1840-JS-20-25-JAN-1840}
			% \label{EMS-1840-JS-DAY-3LETTERMONTH-YEAR}
			
			\begin{description}
				\item[Print Sources]
			\end{description}

			\begin{tabular}{p{0.5in}p{1in}p{0.5in}p{1in}}
				JSP	 	& ---		& JSR 		& --- \\
				DC	 	& ---		& HC 		& --- \\
				HDDC 	& ---		& TCHC 		& --- \\
			\end{tabular}
			% HDDC = woodford's DC diss; JSR = Marquardt's JS Revelations; TCHC = Vogel HC
			
			\begin{description}
				\item[Online Access] \href{https://www.josephsmithpapers.org/paper-summary/letter-to-emma-smith-20-25-january-1840/1}{Here}
				% \item[Analysis] \hyperref[XX]{Here}
			\end{description}
			
			\begin{description}
				\item[Background]
			\end{description}
			
				% It might also be useful to give a two sentence context of the period: e.g., "This speech was given in Nauvoo to a small congregation one month prior to the destruction of the Nauvoo Expositor. These and other tensions may have been on the prophet's mind when he gave the speech."
			
			\begin{description}
				\item[Original Source]
			\end{description}
			
				% brief description of original source material, with reference to JSP or somewhere else for more info
				% Background material: it might be helpful to have a note that says whether a given document was presented as a speech, was written in a journal, etc.
				% Also a comment here about how reliable the source might be, or what shape it is in, if there are large omissions or parts that are hard to understand, and so on.
			
			\begin{description}
				\item[Text]
			\end{description}
			
				\textbf{\sout{Pheadelpha} <Chester Co Pa> [January] 20\supers{th} 1840}

				\textbf{My Dear and beloved Wife}

				\textbf{I recieved a letter from Hyram [Hyrum Smith] which cheared my heart to learn that my Family was all alive yet my heart mourns for those who have been taken from us but not without hope for I shall see them again and be with them therefore we can be more reconciled to the dealings of God I am now makeing all hast[e] to arange my business to start for home I feel very anxious to see you all once more in this world the time seems long that I am deprived of your sosiety but the <lord> being my helper I will not be much longer I am determined to st[art] for home in a few dayes our [bus]iness I expect is before the house of congress now <and> I shall \sout{shall} start for Washington in a few days and from there home as soon as posible I am filled with constant anxiety and shall be until I git home I pray God to spare you all untill I git home my dear Emma my heart is intwined arround you and those little ones I want <you> to remember me tell all \sout{the} chi[l]dren that I love them and will come home as soon as soon as I can yours in the bonds of love your Husband utill [until] Death \&c---}

				\begin{tightcenter}
					\textbf{Joseph Smith Jr}
				\end{tightcenter}

				\textbf{Emma Smith}

				\medskip

				\begin{tightright}
					<25>
				\end{tightright}

				\textbf{Emma Smith}

				\textbf{Commerce}

				\textbf{Hancock \sout{Co}}

				\textbf{County}

				\textbf{Ill---}

				<$\diamond\diamond\diamond$ Han

				Jany 29\supersu{th}\supers{.}>
				
		% -----
		\subsubsection{Discourse, 5 February 1840}
		% -----
			\label{EMS-1840-JS-5-FEB-1840}
			% \label{EMS-1840-JS-DAY-3LETTERMONTH-YEAR}
			
			\begin{description}
				\item[Print Sources]
			\end{description}

			\begin{tabular}{p{0.5in}p{1in}p{0.5in}p{1in}}
				JSP	 	& ---		& JSR 		& --- \\
				DC	 	& ---		& HC 		& --- \\
				HDDC 	& ---		& TCHC 		& --- \\
			\end{tabular}
			% HDDC = woodford's DC diss; JSR = Marquardt's JS Revelations; TCHC = Vogel HC
			
			\begin{description}
				\item[Online Access] \href{https://www.josephsmithpapers.org/paper-summary/discourse-5-february-1840/1}{Here}
				% \item[Analysis] \hyperref[XX]{Here}
			\end{description}
			
			\begin{description}
				\item[Background]
			\end{description}
			
				% It might also be useful to give a two sentence context of the period: e.g., "This speech was given in Nauvoo to a small congregation one month prior to the destruction of the Nauvoo Expositor. These and other tensions may have been on the prophet's mind when he gave the speech."
			
			\begin{description}
				\item[Original Source]
			\end{description}
			
				(Note that I have added the first and last paragraphs, which are not in the JSP entry.)
			
				% brief description of original source material, with reference to JSP or somewhere else for more info
				% Background material: it might be helpful to have a note that says whether a given document was presented as a speech, was written in a journal, etc.
				% Also a comment here about how reliable the source might be, or what shape it is in, if there are large omissions or parts that are hard to understand, and so on.
			
			\begin{description}
				\item[Text]
			\end{description}
			
					\hypertarget{EMS-1840-JS-5-FEB-1840-a}%
				I
					\marginnote{a}%
				went last Evening, to hear \uline{Joe Smith}, the celebrated Mormon, expound his doctrine. I, with several others, had a desire to understand his tenets, as explained by himself. He is not an educated man; but he is a plain, sensible, strong minded man. Every thing he says is said in a manner to leave an impression that he is sincere. There is no le\uline{vity}---no fanaticism---no want of dignity in his deportment. He is apparently from 40 to 45 years of age; rather above the middle stature, and what you ladies would call a good looking man. In his garb there are no peculiarities; his dress being that of a plain, unpretending citizen. He is by profession, a farmer; but is evidently well read.
			
					\hypertarget{EMS-1840-JS-5-FEB-1840-b}%
				He
					\marginnote{b}%
				commenced, by saying, that he knew the prejudices which were abroad in the world against him; but requested us to pay \sout{to} <no> respect to the rumors which were in circulation respecting him and his doctrines. He was accompanied by three or four of his followers. He said--- ``I will state to you our belief, so far as time will permit.''

				I believe, said he, that there is a God, possessing all the attributes ascribed to him by Christians of all denominations. That he reigns over all things in Heaven and on Earth; and that all are subject to his power. He then spoke, rationally, of the attributes of Divinity, such as foreknowledge; mercy, \&c \&c

					\hypertarget{EMS-1840-JS-5-FEB-1840-c}%
				He
					\marginnote{c}%
				then took up the Bible. I believe, said he, in this sacred volume--- In it the Mormon faith is to be found. We teach nothing but what the Bible teaches. We believe nothing but what is to be found in this Book

				I believe in the fall of man, as recorded in the Bible. I believe that God fore-knew every thing; but did not fore-ordain every <thing>. I deny that fore-ordain and fore-know is the same thing. He fore-ordained the fall of man: But all merciful as he is, he fore-ordained, at the same time, a plan of redemption for all mankind. I believe in the Divinity of Jesus Christ, and that he died for the sins of all men, who in Adam had fallen

					\hypertarget{EMS-1840-JS-5-FEB-1840-d}%
				He
					\marginnote{d}%
				then entered into some details, the result of which tended to shew his total unbelief of what is termed \uline{Original} \uline{Sin}. He believes that it is washed away by the blood of Christ, and that it no longer Exists. As a necessary consequence, he believes that we are all born pure and undefiled. That \uline{all} children dying at an Early [age] (say \uline{eight} years) not knowing good from ill, \sout{from} were incapable of sinning; and that all such assuredly go to Heaven.

					\hypertarget{EMS-1840-JS-5-FEB-1840-e}%
				I
					\marginnote{e}%
				believe, said he, that a man is a moral, responsible, free agent, that although it was fore-ordained he should fall and be redeemed; yet after the redemption it was not fore-ordained that he should again sin. In the Bible a rule of conduct is laid down for <him>. In the Old and New testaments the law by which he is to be governed may be found. If he violates that law, he is to be punished for the deeds done in the body

					\hypertarget{EMS-1840-JS-5-FEB-1840-f}%
				I
					\marginnote{f}%
				believe that God is Eternal. That he had no beginning and can have no End. Eternity means that which is without beginning or End. I believe that the \uline{Soul} is Eternal. It had no beginning; it can have no End. Here he entered into some Explanations, which were so brief that I could not perfectly comprehend him. But the idea seemed to be that the Soul of man--- the Spirit, had Existed from Eternity in the bosom of Divinity, and so far as he was intelligible to me, must ultimately return from whence it came

					\hypertarget{EMS-1840-JS-5-FEB-1840-g}%
				He
					\marginnote{g}%
				said very little of rewards and punishments. But one conclusion, from what he did say, was irresistible. He contended throughout, that every thing which had a \uline{beginning} must have an \uline{Ending}; and consequently, if the punishment of man \uline{commences} in the next world, it must, according to his logic and belief have an \uline{End}

					\hypertarget{EMS-1840-JS-5-FEB-1840-h}%
				During
					\marginnote{h}%
				the whole of his address, and it occupied more than two hours, there was no opinion or belief that he Expressed, that was calculated, in the slightest degree, to impair the morals of Society, or in any manner to degrade and brutalize the human species. There was much in his precepts, if they were followed, that would soften the asperities of man towards man, and that would tend to make him a more rational being than he is generally found to be. There was no violence; no fury; no denunciation. His religion appears to be the religion of meekness; lowliness, and mild persuasion.

					\hypertarget{EMS-1840-JS-5-FEB-1840-i}%
				\phantom{ }
					\marginnote{i}%
				\sout{To} <Towards> the close of his address, he remarked that he had been represented as pretending to be a Saviour, a Worker of Miracles, \&c. All this was false. He made no such pretensions. He was but a man, he said--- a plain untutored <man>; seeking what he should do to be saved. He performed no miracles. He did not pretend to possess any such power

					\hypertarget{EMS-1840-JS-5-FEB-1840-j}%
				He
					\marginnote{j}%
				closed by referring to the Mormon Bible, which, he said, contained nothing inconsistent or conflicting with the Christian Bible. And he again repeated,

				\begin{quote}
					that all who would follow the precepts of the Bible whether Mormon or not, would assuredly be saved. Throughout his whole address, he displayed, strongly, a spirit of Charity and Forbearance
				\end{quote}

				The Mormon Bible, he said, was communicated to him, \uline{direct from Heaven}. If there was such a thing on Earth, as the Author of it, then he (Smith) was the Author; but the idea that he wished to impress was, that he had penned it as dictated by God.
				
					\hypertarget{EMS-1840-JS-5-FEB-1840-k}%
				I
					\marginnote{k}%
				have taken some pains to explain this man's belief, as he himself explained it. I have done so, because, I thought it might satisfy your curiosity, \& might be interesting to you and some of your friends. I have changed my opinion of the Mormons. They are an injured and much abused people. Of matters of \uline{Faith}, you know I express no opinion.
				
				I have only room to add---let William, if you cannot do it acknowledge the receipt of this, with the inclosure. Remember me to Sarah and the Boys. Kiss the dear Baby for me. Affectionately, your Husband
				
				\begin{tightright}
					\uline{M. L. D}avis
				\end{tightright}
				
					\hypertarget{EMS-1840-JS-5-FEB-1840-l}%
				I
					\marginnote{l}%
				omitted to say, he does not believe in infant baptism, \uline{sprinkling} but in \uline{immersion}, after \uline{eight} years of age
				
		% -----
		\subsubsection{Discourse, 1 March 1840}
		% -----
			\label{EMS-1840-JS-1-MAR-1840}
			% \label{EMS-1840-JS-DAY-3LETTERMONTH-YEAR}
			
			\begin{description}
				\item[Print Sources]
			\end{description}

			\begin{tabular}{p{0.5in}p{1in}p{0.5in}p{1in}}
				JSP	 	& ---		& JSR 		& --- \\
				DC	 	& ---		& HC 		& --- \\
				HDDC 	& ---		& TCHC 		& --- \\
			\end{tabular}
			% HDDC = woodford's DC diss; JSR = Marquardt's JS Revelations; TCHC = Vogel HC
			
			\begin{description}
				\item[Online Access] \href{https://www.josephsmithpapers.org/paper-summary/discourse-1-march-1840/1}{Here}
				% \item[Analysis] \hyperref[XX]{Here}
			\end{description}
			
			\begin{description}
				\item[Background]
			\end{description}
			
				% It might also be useful to give a two sentence context of the period: e.g., "This speech was given in Nauvoo to a small congregation one month prior to the destruction of the Nauvoo Expositor. These and other tensions may have been on the prophet's mind when he gave the speech."
			
			\begin{description}
				\item[Original Source]
			\end{description}
			
				% brief description of original source material, with reference to JSP or somewhere else for more info
				% Background material: it might be helpful to have a note that says whether a given document was presented as a speech, was written in a journal, etc.
				% Also a comment here about how reliable the source might be, or what shape it is in, if there are large omissions or parts that are hard to understand, and so on.
			
			\begin{description}
				\item[Text]
			\end{description}
			
					\hypertarget{EMS-1840-JS-1-MAR-1840-a}%
				After
					\marginnote{a}%
				engaging in prayer to the Most High, and reading a chapter of sacred writ, he commenced his discourse. He told his people he was their servant; that they had a right to know all the incidents of his journey; he would therefore endeavor to give them a minute account. He did not like to preach politics on the Sabbath; but he must free his mind,---must tell the whole story.

					\hypertarget{EMS-1840-JS-1-MAR-1840-b}%
				The
					\marginnote{b}%
				object of his visit at Washington, you well know, was to make application to Congress for relief, touching their troubles in Missouri. But to the discourse. He said, on his arrival in Washington, he, with two of his Elders, ([Sidney] Rigdon and [Elias] Higbee,) called on Mr. [Martin] Van Buren at the ``White House'' with a letter of introduction, and after making known to him the object of their visit, and soliciting him to help them, Mr. Van Buren replied; ``Help you! how can I help you? All Missouri would turn against me.'' But they demanded of Mr. Van Buren a hearing, and he, after listening a few moments to their tale of injured innocence, abruptly left the room. After waiting some time for his return, they were under the necessity of departing, disappointed, and chagrined.

					\hypertarget{EMS-1840-JS-1-MAR-1840-c}%
				He
					\marginnote{c}%
				thought Mr. Van Buren treated them with great disrespect and neglect. He said while they were with Mr. Van Buren, a member of Congress waited upon him, and in conversation, among other things, told the President that he (the President) was getting fat. The President replied that he was aware of the fact; that he had to go every few days to the tailor's to get his clothes let out, or purchase a new coat. The ``Prophet'' here added, at the top of his voice,--- \textit{he hoped he would continue to grow fat, and swell, and, before the next election, burst!}

					\hypertarget{EMS-1840-JS-1-MAR-1840-d}%
				He
					\marginnote{d}%
				felt at home in the White House, and, while there, thought he began to swell a little himself. He felt that he had a perfect right there, as much right as Van Buren, because it belonged to the people, and he was one of the people.

				He spoke of the success attending his preaching in the Eastern cities; that the people \textit{en masse} in many places were converted to the Mormon faith. That the striplings which they had sent from this wilderness to preach to the wise men and priests of the great cities, were accomplishing mighty things for the church, by confounding the learned priests.

					\hypertarget{EMS-1840-JS-1-MAR-1840-e}%
				On
					\marginnote{e}%
				one occasion, he said, six ministers attended a meeting where a little Mormon fellow was preaching, and undertook to put him down by ridicule: but he stood his ground, and whipped them all out, by fair argument; and the congergation acknowledged \textit{en masse} that he whipped them all.

				Here this \textit{spiritual} discourse was brought to a close, by a violent shower of rain. After making an appointment to deliver the conclusion, the Prophet dismissed the meeting.
				
		% -----
		\subsubsection{Minutes and Discourse, 6 March 1840}
		% -----
			\label{EMS-1840-JS-6-MAR-1840}
			% \label{EMS-1840-JS-DAY-3LETTERMONTH-YEAR}
			
			% --
			\paragraph{As Reported by Elias Smith}
			% --
				\label{EMS-1840-JS-6-MAR-1840-ES}
			
				\begin{description}
					\item[Print Sources]
				\end{description}

				\begin{tabular}{p{0.5in}p{1in}p{0.5in}p{1in}}
					JSP	 	& ---		& JSR 		& --- \\
					DC	 	& ---		& HC 		& --- \\
					HDDC 	& ---		& TCHC 		& --- \\
				\end{tabular}
				% HDDC = woodford's DC diss; JSR = Marquardt's JS Revelations; TCHC = Vogel HC
				
				\begin{description}
					\item[Online Access] \href{https://www.josephsmithpapers.org/paper-summary/minutes-and-discourse-6-march-1840/1}{Here}
					% \item[Analysis] \hyperref[XX]{Here}
				\end{description}
				
				\begin{description}
					\item[Background]
				\end{description}
				
					% It might also be useful to give a two sentence context of the period: e.g., "This speech was given in Nauvoo to a small congregation one month prior to the destruction of the Nauvoo Expositor. These and other tensions may have been on the prophet's mind when he gave the speech."
				
				\begin{description}
					\item[Original Source]
				\end{description}
				
					% brief description of original source material, with reference to JSP or somewhere else for more info
					% Background material: it might be helpful to have a note that says whether a given document was presented as a speech, was written in a journal, etc.
					% Also a comment here about how reliable the source might be, or what shape it is in, if there are large omissions or parts that are hard to understand, and so on.
				
				\begin{description}
					\item[Text]
				\end{description}
				
						\hypertarget{EMS-1840-JS-6-MAR-1840-ES-a}%
					...President
						\marginnote{a}%
					Joseph Smith Jun. then addressed the Council on various subjects, \&, in particular, the Consecration Law; Stating, that the affair now before Congress was the only thing that ought to interest the saints at present. \& till it was ascertained how it would terminate no person ought to be brought to account before the constituted authorities of the Church for any offence whatever, \& was determined that no man should be brought before the Council in Nauvoo till that time \&c. \&c.
					
						\hypertarget{EMS-1840-JS-6-MAR-1840-ES-b}%
					He
						\marginnote{b}%
					said that the Law of consecration could not be kept here, \& that it was the will of the Lord that we should desist from trying to keep it, \& if persisted in it would produce a perfect abortion, \& that he assumed the whole responsibility of not keeping it untill proposed by himself. He requested every exertion to be made to forward affidavits to Washington, \& also letters to members of Congress.
					
					The following votes were then passed.
					
					1st. that this Council will coincide with Pr. Joseph Smith jr. decission concerning the consecration Law, on the principle of its being the will of the Lord, \& of Pr. Smith's taking the responsibility on himself...
				
			% --
			\paragraph{As Reported by John Smith}
			% --
				\label{EMS-1840-JS-6-MAR-1840-JS}
			
				\begin{description}
					\item[Print Sources]
				\end{description}

				\begin{tabular}{p{0.5in}p{1in}p{0.5in}p{1in}}
					JSP	 	& ---		& JSR 		& --- \\
					DC	 	& ---		& HC 		& --- \\
					HDDC 	& ---		& TCHC 		& --- \\
				\end{tabular}
				% HDDC = woodford's DC diss; JSR = Marquardt's JS Revelations; TCHC = Vogel HC
				
				\begin{description}
					\item[Online Access] \href{https://www.josephsmithpapers.org/paper-summary/discourse-6-march-1840-as-reported-by-john-smith/1}{Here}
					% \item[Analysis] \hyperref[XX]{Here}
				\end{description}
				
				\begin{description}
					\item[Background]
				\end{description}
				
					% It might also be useful to give a two sentence context of the period: e.g., "This speech was given in Nauvoo to a small congregation one month prior to the destruction of the Nauvoo Expositor. These and other tensions may have been on the prophet's mind when he gave the speech."
				
				\begin{description}
					\item[Original Source]
				\end{description}
				
					% brief description of original source material, with reference to JSP or somewhere else for more info
					% Background material: it might be helpful to have a note that says whether a given document was presented as a speech, was written in a journal, etc.
					% Also a comment here about how reliable the source might be, or what shape it is in, if there are large omissions or parts that are hard to understand, and so on.
				
				\begin{description}
					\item[Text]
				\end{description}
				
					March 6th \sout{1843} 1840
					
					Met to in Council Joseph \& Hyrum [Smith] present who informed us that thus Saith the <Lord> you need not observe the Law of Consecration untill our case was Decided in Congress also to be very mild in our Difficulties with each other not have any trials in council if it could be avoided would to Lord that we could be at peace among ourselves
					
		% -----
		\subsubsection{Letter to Elias Higbee, 7 March 1840}
		% -----
			\label{EMS-1840-JS-7-MAR-1840}
			% \label{EMS-1840-JS-DAY-3LETTERMONTH-YEAR}
			
			\begin{description}
				\item[Print Sources]
			\end{description}

			\begin{tabular}{p{0.5in}p{1in}p{0.5in}p{1in}}
				JSP	 	& ---		& JSR 		& --- \\
				DC	 	& ---		& HC 		& --- \\
				HDDC 	& ---		& TCHC 		& --- \\
			\end{tabular}
			% HDDC = woodford's DC diss; JSR = Marquardt's JS Revelations; TCHC = Vogel HC
			
			\begin{description}
				\item[Online Access] \href{https://www.josephsmithpapers.org/paper-summary/letter-to-elias-higbee-7-march-1840/1}{Here}
				% \item[Analysis] \hyperref[XX]{Here}
			\end{description}
			
			\begin{description}
				\item[Background]
			\end{description}
			
				% It might also be useful to give a two sentence context of the period: e.g., "This speech was given in Nauvoo to a small congregation one month prior to the destruction of the Nauvoo Expositor. These and other tensions may have been on the prophet's mind when he gave the speech."
			
			\begin{description}
				\item[Original Source]
			\end{description}
			
				Letter written by Elias Smith, with postscript of approval by JS.
			
				% brief description of original source material, with reference to JSP or somewhere else for more info
				% Background material: it might be helpful to have a note that says whether a given document was presented as a speech, was written in a journal, etc.
				% Also a comment here about how reliable the source might be, or what shape it is in, if there are large omissions or parts that are hard to understand, and so on.
			
			\begin{description}
				\item[Text]
			\end{description}
			
					\hypertarget{EMS-1840-JS-7-MAR-1840-a}%
				...I
					\marginnote{a}%
				am instructed by the unanimous voice of the Church here to request you not to give an inch in any position that has been taken; especially in relation to our damages, the amount of which we think has been estimated full low, and if an offer should be made to remunerate us for our Lands at entry price, or any thing short of the full amount claimed in any particular--- to treat such offer with proper contempt. Our rights we ask for; our rights we want, and must have of the Government of the United States: And nothing short of that can be received.
					\hypertarget{EMS-1840-JS-7-MAR-1840-b}%
				We
					\marginnote{b}%
				demand them upon constitutional premises, and expect to [be] heard in our appeal to the General Government, and answered in a manner which shall be honorable to the guardians, of the rights and privileges of American citizens; who compose that honorable body, and acceptable to those who have been persecuted, Smitten, and driven from their peaceful habitations, in violation of the Constitution, and of the natural and unalienable rights of man---

					\hypertarget{EMS-1840-JS-7-MAR-1840-c}%
				Our
					\marginnote{c}%
				cause will unquestionably be strongly opposed by such, relentless, unfeeling and unprincipled demagogues as, [Thomas Hart] Benton, and others of the ``Golden humbug'', firm of Vanburen [Martin Van Buren] \&, co. who have disgraced the American character, by their subserviency to a s[c]heme of [s]peculations and political fraud, The sole aim of whom has been, to aggrandize themselves, at the expence of the labouring class of community, and to enrich the few at expence of the many. As regardless of honor, honesty; or the rights of man, as Bonapart, Nero, or Caligula--- I cincerely hope that every such unprincipled character, will take a decided stand in favor of Bogg's [Lilburn W. Boggs's] exterminating decree; as such a course would speak volumes in our favor, and expose them to the just indignation of an abused and injured publick, who will not fail to meet \sout{ours} out to them a just recompence of reward without distinctions of name, sect or party by applying that peaceable and constitutional corrective, guarenteed to every American citizen, the right of suffrage, if there is virtue enough left in the sons of the revolutionary Fathers, to transmit to their posterity, those inestimable blessings purchased by Fathers at the expense of so much blood and Treasure.
					\hypertarget{EMS-1840-JS-7-MAR-1840-d}%
				With
					\marginnote{d}%
				these feelings, and desires, and believing that, that Being, who holds the destinies of nations of the earth in his hands, will overrule all things for the furtherance of the work, which he has set his hands to perform in these days, till the throne shall be cast down, The Ancient of days sit and the Kingdom and dominion, and the greatness of the Kingdom under the whole heavens shall be given to the people of the Saints of the most high--- Whose kingdom is an everlasting Kingdom, and dominions shall serve and obey him,

				I subscribe myself Yours \&c

				\begin{tightright}
					Elias Smith
				\end{tightright}

				Hon E[lias] Higbee

				P.S.

				Dear Brother

				I heartily concur in the above sentiment

				\begin{tightcenter}
					Yours with sentiments
				\end{tightcenter}

				\begin{tightright}
					And feelings of the \sout{kind} best kind

					Josiph Smith J\supers{r.}
				\end{tightright}
				
		% -----
		\subsubsection{Letter to Robert D. Foster, 11 March 1840}
		% -----
			\label{EMS-1840-JS-11-MAR-1840}
			% \label{EMS-1840-JS-DAY-3LETTERMONTH-YEAR}
			
			\begin{description}
				\item[Print Sources]
			\end{description}

			\begin{tabular}{p{0.5in}p{1in}p{0.5in}p{1in}}
				JSP	 	& ---		& JSR 		& --- \\
				DC	 	& ---		& HC 		& --- \\
				HDDC 	& ---		& TCHC 		& --- \\
			\end{tabular}
			% HDDC = woodford's DC diss; JSR = Marquardt's JS Revelations; TCHC = Vogel HC
			
			\begin{description}
				\item[Online Access] \href{https://www.josephsmithpapers.org/paper-summary/letter-to-robert-d-foster-11-march-1840/1}{Here}
				% \item[Analysis] \hyperref[XX]{Here}
			\end{description}
			
			\begin{description}
				\item[Background]
			\end{description}
			
				See \href{https://www.josephsmithpapers.org/paper-summary/minutes-20-december-1840/1}{Minutes, 20 December 1840}
			
				% It might also be useful to give a two sentence context of the period: e.g., "This speech was given in Nauvoo to a small congregation one month prior to the destruction of the Nauvoo Expositor. These and other tensions may have been on the prophet's mind when he gave the speech."
			
			\begin{description}
				\item[Original Source]
			\end{description}
			
				% brief description of original source material, with reference to JSP or somewhere else for more info
				% Background material: it might be helpful to have a note that says whether a given document was presented as a speech, was written in a journal, etc.
				% Also a comment here about how reliable the source might be, or what shape it is in, if there are large omissions or parts that are hard to understand, and so on.
			
			\begin{description}
				\item[Text]
			\end{description}
			
				\begin{tightright}
						\hypertarget{EMS-1840-JS-11-MAR-1840-a}%
					Nauvoo
						\marginnote{a}%
					March 11, 1840.
				\end{tightright}
				
				Sir,

				After I left you, I came to my bro's house in Plymouth the same day; and there I learned that my father was sick, and that he was not expected to live--- had called his children together \&c. my Bro. William [Smith] had left home for this place the day before I arrived there. But when I arrived here I found him a little better, but was quite low yet. Since that time, he has been much afflicted with the ague, but is now recovering. With that exception we are all well at present; and it is a general time of health here now.

					\hypertarget{EMS-1840-JS-11-MAR-1840-b}%
				I
					\marginnote{b}%
				have delivered two discourses in this place since my return---giveing a brief history of our journy the reception we met with by the president \&c. and the general feeling towards us in Washington and other places.%
					\footnote{%
						% \label{}%
						JSP note: ``JS delivered one of these two discourses in the Commerce area on 1 March 1840. In that discourse, JS stated he felt that President Martin Van Buren had `treated them with great disrespect and neglect' and that missionaries in the eastern United States were having great success. It is unknown when JS delivered the other discourse. In letters home in late 1839, JS reported that the church's delegates to the federal government were well received by many in Congress and that Washington residents listened to them with great interest. (Discourse, 1 Mar. 1840; Letter to Hyrum Smith and Nauvoo High Council, 5 Dec. 1839; Letter to Seymour Brunson and Nauvoo High Council, 7 Dec. 1839.) ''
						}
				The effect has been to turn the entire mass of the people, even to an individual, so far as I have learned on the other side of the great political question---

				I find that we have lost nothing by our change; but have gained friends and influence. The fact is, we were compelled to change in consequence of seeing a disposition manifest to turn a deaf ear to the cries of suffering innocence. When we can see a disposition in our chief magistrate to sacrifice the rights of the poor at the shrine of popularity, it is high time to cast off such an individual.

					\hypertarget{EMS-1840-JS-11-MAR-1840-c}%
				After
					\marginnote{c}%
				haveing formed an acquaintance with you, and a very intimate one too, for the last 4 months, and I need not say an agreeable one \sout{too}, I feel quite anxious to see you after a short separation, I hope you can make it convenient to come up and see us soon. I want to get hold of your journal very much.%
					\footnote{%
						% \label{}%
						JSP note: ``JS either had tasked Foster with keeping his journal during their travels throughout the eastern United States or desired to access Foster's journal in order to record the details of the trip in his own personal record. Years later, the compilers of the manuscript history of the church reported JS declaring, `I depended on Dr Foster to keep my daily journal during this journey but he has failed me.' (Historian's Office, JS History, Draft Notes, 4 Mar. 1840, 5.)''
						}

				Our Church here is prospering, and many are comeing into it. Our Town is improveing very fast. It is almost incredible to see what amt. of labor has been performed here during the winter past. There is now every prospect of our haveing a good society, a peaceable habitation and a desirable residence here.

				May the Lord prosper our righteous cause, and save us in the day of his comeing!

				\begin{tightcenter}
					As ever,
				\end{tightcenter}

				\begin{tightright}
					I am Your friend And Brother in the New and evelasting Cov't.
					
					Joseph Smith Jun.
				\end{tightright}

					\hypertarget{EMS-1840-JS-11-MAR-1840-d}%
				Robt.
					\marginnote{d}%
				D Foster M. D.

				Beverley

				Adams Co

				Illinois

				P.S. Our Business in Washington has gone before \sout{the} <a> committee on <the> judiciary without a dissenting voice. We have recently received two letters from Judge [Elias] Higbee to this effect. He is well, But pres't. [Sidney] Rigdon yet has the chills and fever.

				\begin{tightright}
					Our best respects to Bro. Wilbur \& family; and to all other friends in that section
					
					\uline{As before}

					\uline{J.S. Jun}
				\end{tightright}


				[\textit{page [3] blank}]

				\begin{tightright}
					<11>
				\end{tightright}

				\begin{tightcenter}
					Robert D. Foster M. D.
				\end{tightcenter}

				\begin{tightright}
					Beverley Post Office

					Adams County

					Illinois
				\end{tightright}

				Commerce Ills
				
				March 13
				
		% -----
		\subsubsection{Minutes and Discourse, 6--8 April 1840}
		% -----
			\label{EMS-1840-JS-6-8-APR-1840}
			% \label{EMS-1840-JS-DAY-3LETTERMONTH-YEAR}
			
			\begin{description}
				\item[Print Sources]
			\end{description}

			\begin{tabular}{p{0.5in}p{1in}p{0.5in}p{1in}}
				JSP	 	& ---		& JSR 		& --- \\
				DC	 	& ---		& HC 		& --- \\
				HDDC 	& ---		& TCHC 		& --- \\
			\end{tabular}
			% HDDC = woodford's DC diss; JSR = Marquardt's JS Revelations; TCHC = Vogel HC
			
			\begin{description}
				\item[Online Access] \href{https://www.josephsmithpapers.org/paper-summary/minutes-and-discourse-6-8-april-1840/1}{Here}
				% \item[Analysis] \hyperref[XX]{Here}
			\end{description}
			
			\begin{description}
				\item[Background]
			\end{description}
			
				% It might also be useful to give a two sentence context of the period: e.g., "This speech was given in Nauvoo to a small congregation one month prior to the destruction of the Nauvoo Expositor. These and other tensions may have been on the prophet's mind when he gave the speech."
			
			\begin{description}
				\item[Original Source]
			\end{description}
			
				% brief description of original source material, with reference to JSP or somewhere else for more info
				% Background material: it might be helpful to have a note that says whether a given document was presented as a speech, was written in a journal, etc.
				% Also a comment here about how reliable the source might be, or what shape it is in, if there are large omissions or parts that are hard to understand, and so on.
			
			\begin{description}
				\item[Text]
			\end{description}
			
					\hypertarget{EMS-1840-JS-6-8-APR-1840-a}%
				...The
					\marginnote{a}%
				president rose and made some observations on the business of the conference; exhorted the brethren who had charges to bring against any individual to be charitable; and made some very appropriate remarks respecting ``pulling out the beam in their own eyes, that they might see clearly the mote which was in their brothers eye...
				
				The President called upon the Clerk to read the report of the Presidency and High council, with regard to their proceedings in purchasing lands and securing a place of gathering for the saints. The report having been read, the President made some observations respecting the pecuniary affairs of the church, and requested the brethren to step forward and assist in liquidating the debts on the town plot, so that the poor might have inheritances.
				
					\hypertarget{EMS-1840-JS-6-8-APR-1840-b}%
				He
					\marginnote{b}%
				then gave some account of his mission to Washington city, in company with President [Sidney] Rigdon and Judge [Elias] Higbee, the treatment they received and the action of the Senate on the memorial which was presented to them...
				
				Conference met pursuant to adjournment, after singing the President arose and read the 3d chap. of John's Gospel after which prayer was offered by elder Erastus Snow.
				
				The President commenced making observations on the different subjects embraced in the chapter particularly on the 3d, 4th, 5th verses illustrating it with a very beautiful and striking figure, and throwing a flood of light on the subjects which were brought up to review.
				
					\hypertarget{EMS-1840-JS-6-8-APR-1840-c}%
				He
					\marginnote{c}%
				then spoke to the elders respecting their mission, and advised those who went into the world, to preach the gospel, to leave their families provided for, with the necessaries of life; and to teach the gathering as set forth in the Holy scriptures.
				
				That it had been wisdom to, for the greater body of the church to keep on this side of the river, in order that a foundation might be established in this place, but that now, it was the priviledge of the saints to occupy the lands in the Iowa, or wherever the spirit might lead them.
				
				That he did not wish to have any political influence, but wished the saints to use their political franchise to the best of their knowledge.
				
				He then stated that since Elder Hyde had been appointed to visit the Jewish people, he had felt an impression that it would be well for Elder John E. Page to accompany him on his mission...
				
		% -----
		\subsubsection{Discourse, 7 April 1840}
		% -----
			\label{EMS-1840-JS-7-APR-1840}
			% \label{EMS-1840-JS-DAY-3LETTERMONTH-YEAR}
			
			\begin{description}
				\item[Print Sources]
			\end{description}

			\begin{tabular}{p{0.5in}p{1in}p{0.5in}p{1in}}
				JSP	 	& ---		& JSR 		& --- \\
				DC	 	& ---		& HC 		& --- \\
				HDDC 	& ---		& TCHC 		& --- \\
			\end{tabular}
			% HDDC = woodford's DC diss; JSR = Marquardt's JS Revelations; TCHC = Vogel HC
			
			\begin{description}
				\item[Online Access] \href{https://www.josephsmithpapers.org/paper-summary/discourse-7-april-1840/1}{Here}
				% \item[Analysis] \hyperref[XX]{Here}
			\end{description}
			
			\begin{description}
				\item[Background]
			\end{description}
			
				% It might also be useful to give a two sentence context of the period: e.g., "This speech was given in Nauvoo to a small congregation one month prior to the destruction of the Nauvoo Expositor. These and other tensions may have been on the prophet's mind when he gave the speech."
			
			\begin{description}
				\item[Original Source]
			\end{description}
			
				% brief description of original source material, with reference to JSP or somewhere else for more info
				% Background material: it might be helpful to have a note that says whether a given document was presented as a speech, was written in a journal, etc.
				% Also a comment here about how reliable the source might be, or what shape it is in, if there are large omissions or parts that are hard to understand, and so on.
			
			\begin{description}
				\item[Text]
			\end{description}
			
					\hypertarget{EMS-1840-JS-7-APR-1840-a}%
				He
					\marginnote{a}%
				said that soon after reaching Washington, he called on Mr. [Martin] Van Buren, and asked permission to leave with him the memorial with which he had been entrusted, at the same time briefly stating its contents. Mr. Van Buren's manner was very repulsive, and it was only after his (Smith's) urgent request that he consented to receive the paper and to give an answer on the morrow. The next day Smith again called, when Mr. Van Buren cut short the interview by saying, ``I can do nothing for you, gentlemen. If I were to, I should go against the whole state of Missouri, and that state would go against me at the next election.''
					\hypertarget{EMS-1840-JS-7-APR-1840-b}%
				Mr.
					\marginnote{b}%
				Smith said he was thunderstruck at this avowal. He had always believed Mr. Van Buren to be a high-minded statesman, and had uniformly supported him as such; but he now saw that he was only a huckstering politician, who would sacrifice any and every thing to promote his re-election. He left him abruptly, and rejoiced, when without the walls of the palace, that he could once more breathe the air of a freeman.
				
		% -----
		\subsubsection{Minutes, 12 April 1840}
		% -----
			\label{EMS-1840-JS-12-APR-1840}
			% \label{EMS-1840-JS-DAY-3LETTERMONTH-YEAR}
			
			\begin{description}
				\item[Print Sources]
			\end{description}

			\begin{tabular}{p{0.5in}p{1in}p{0.5in}p{1in}}
				JSP	 	& ---		& JSR 		& --- \\
				DC	 	& ---		& HC 		& --- \\
				HDDC 	& ---		& TCHC 		& --- \\
			\end{tabular}
			% HDDC = woodford's DC diss; JSR = Marquardt's JS Revelations; TCHC = Vogel HC
			
			\begin{description}
				\item[Online Access] \href{https://www.josephsmithpapers.org/paper-summary/minutes-12-april-1840/1}{Here}
				% \item[Analysis] \hyperref[XX]{Here}
			\end{description}
			
			\begin{description}
				\item[Background]
			\end{description}
			
				% It might also be useful to give a two sentence context of the period: e.g., "This speech was given in Nauvoo to a small congregation one month prior to the destruction of the Nauvoo Expositor. These and other tensions may have been on the prophet's mind when he gave the speech."
			
			\begin{description}
				\item[Original Source]
			\end{description}
			
				% brief description of original source material, with reference to JSP or somewhere else for more info
				% Background material: it might be helpful to have a note that says whether a given document was presented as a speech, was written in a journal, etc.
				% Also a comment here about how reliable the source might be, or what shape it is in, if there are large omissions or parts that are hard to understand, and so on.
			
			\begin{description}
				\item[Text]
			\end{description}
			
					\hypertarget{EMS-1840-JS-12-APR-1840-a}%
				April
					\marginnote{a}%
				12\supers{th} 1840 The High Council of the Church of Jesus Christ of Latter day Saints of Nauvoo I\supersu{ll} met in council at the house of J. Smith j\supers{r.}

				1 A charge was prefered against Alonson [Alanson] Ripley by Alva Keller for taking rails from his lot.

				After the matter had been explained by Pres\supers{t} Joseph Smith j\supers{r}, the charge was withdrawn.

				2\supers{nd} President Joseph Smith j\supersu{r} spoke relative to Elder Oliver Gran[g]er's going to the East to settle some buisness transactions for the church, recomending that the Council should appoint some one to go with him.---

					\hypertarget{EMS-1840-JS-12-APR-1840-b}%
				3\supers{rd}
					\marginnote{b}%
				On motion it was voted that Hyrum Smith go with O. Granger to assist him in the aforesaid buisness transaction.

				4\supers{th} Voted that Pres\supersu{t} Joseph Smith j\supersu{r} make the necessary credentials for Oliver Granger and Hyrum Smith, concerning their buisness transactions in the East.---

				Adjourned.

				\begin{tightright}
					Hosea Stout Clerk Pro tem.---
				\end{tightright}

				Recorded on the 14\supers{th} February 1842 by Hosea Stout Clerk
				
		% -----
		\subsubsection{Letter to Orson Hyde and John E. Page, 14 May 1840}
		% -----
			\label{EMS-1840-JS-14-MAY-1840}
			% \label{EMS-1840-JS-DAY-3LETTERMONTH-YEAR}
			
			\begin{description}
				\item[Print Sources]
			\end{description}

			\begin{tabular}{p{0.5in}p{1in}p{0.5in}p{1in}}
				JSP	 	& ---		& JSR 		& --- \\
				DC	 	& ---		& HC 		& --- \\
				HDDC 	& ---		& TCHC 		& --- \\
			\end{tabular}
			% HDDC = woodford's DC diss; JSR = Marquardt's JS Revelations; TCHC = Vogel HC
			
			\begin{description}
				\item[Online Access] \href{https://www.josephsmithpapers.org/paper-summary/letter-to-orson-hyde-and-john-e-page-14-may-1840/1}{Here}
				% \item[Analysis] \hyperref[XX]{Here}
			\end{description}
			
			\begin{description}
				\item[Background]
			\end{description}
			
				% It might also be useful to give a two sentence context of the period: e.g., "This speech was given in Nauvoo to a small congregation one month prior to the destruction of the Nauvoo Expositor. These and other tensions may have been on the prophet's mind when he gave the speech."
			
			\begin{description}
				\item[Original Source]
			\end{description}
			
				% brief description of original source material, with reference to JSP or somewhere else for more info
				% Background material: it might be helpful to have a note that says whether a given document was presented as a speech, was written in a journal, etc.
				% Also a comment here about how reliable the source might be, or what shape it is in, if there are large omissions or parts that are hard to understand, and so on.
			
			\begin{description}
				\item[Text]
			\end{description}
			
				\begin{tightright}
						\hypertarget{EMS-1840-JS-14-MAY-1840-a}%
					Nauvoo
						\marginnote{a}%
					Hancock Co. Ills May 14\supers{th} 1840
				\end{tightright}

				Dear Bretheren

				I am happy in being informed by your letter that your mission swells ``larger and larger''; it is a great and important mission, and one that is worthy of those inteligences%
					\footnote{%
						% \label{}%
						JSP note: ``An 1833 revelation described the eternal nature of `inteligences,' teaching that `man was also in the begining with God, inteligence or the Light of truth was not created or made neith[er] indeed can be.' (Revelation, 6 May 1833 [D\&C 93:29].)''
						}
				who surround the throne of Jehovah to be ingaged in; Altho it appears great at present, yet you have but just begun to realize the greatness, the extent and glory of the same. If there is any thing calculated to interest the mind of the saints, to awaken in them the finest sensibilities; and arouse them to enterprise, and exertion, surely it is the great and precious promises, made by our heavenly Father to the children of Abraham; and those engaged in seeking the outcasts of Israel, and the dispersed of Judah, cannot fail to enjoy the Spirit of the Lord, and have the choisests blessings of Heaven rest upon them in copious effusions, Bretheren you are in the path way to Eternal Fame! and immortal Glory; and inasmuch as you feel interested for the covenant people of the Lord, the God of their Father shall bless you.
					\hypertarget{EMS-1840-JS-14-MAY-1840-b}%
				Do
					\marginnote{b}%
				not be discouraged on accou[n]t of the greatness of the work; only be humble, and faithful, and then you can say, ``what art thou, O, great mountain, ``before Zerubbable shalt thou be brought down'' He who scattered Israel has promised to gather them; therefore, inasmuch as you are to be instrumental in this great work, he will endow you with power, wisdom, might, and inteligence; and every qualification necessary; while your minds will expand wider and wider, untill you can circumscribe the Earth, \& the Heavens, and reach forth into eternity; contemplate the mighty acts of Jehovah, in all their variety \& glory

					\hypertarget{EMS-1840-JS-14-MAY-1840-c}%
				In
					\marginnote{c}%
				answer to your inquiries, respecting the translation and publication, of the Book of Mormon, Hymn Book, History of the church, \&c, \&c; I would say, that I entirely approve of the same; and give my consent, with the exception of the Hymn Book, as a new edition, containing a greater variety of Hymns, will be shortly published or printed in this place; which, I think will be a standard work

				As soon as it is printed, you shall have some to you, which you may get translated, and printed into any language you please. Should we not be able to send some to you, and there should be a great call for Hymns where you \sout{are} may be; then I should have no objections to your publishing the present one. Were you to publish the Book of Mormon, Doctrine \& Covenants, or Hymn Book; I desire the copy rights of the same to be secured in my name.

					\hypertarget{EMS-1840-JS-14-MAY-1840-d}%
				With
					\marginnote{d}%
				respect to publishing any other work, either original, or those which have been published before, you will be governed by circumstances; if you think necessary to do so I shall have no objections whatever--- It will be well to study plainness and simplicity in whatever you may publish ``for my soul delighteth in pla[i]nness''.

				I feel much pleased with the spirit of your letter, and be assured, Dear Bretheren, of my hearty co-operation, and my prayers for your welfare and sucess.

					\hypertarget{EMS-1840-JS-14-MAY-1840-e}%
				In
					\marginnote{e}%
				answer to your enquiry in a former letter, relative to the duty of the seventies in regulating churches \&c; I say that the duties of the seventies is, more particularly to preach the Gospel, \& build up churches, rather than regulate them.--- that a high priest may take charge of them.

				If a high priest should be remiss in his duty, \& should lead, or suffer the church to be led astray; depart from the ordinances of the Lord, then it is the duty of one of the seventies, acting under the special direction of the twelve---being duly commissioned by them with their delegated authority, to go that church and if agreeable to a majority of the members of said church to proceed to regulate and put in order the same--- otherwise he can have no authority to act
				
		% -----
		\subsubsection{Memorial to Nauvoo High Council, 18 June 1840}
		% -----
			\label{EMS-1840-JS-18-JUN-1840}
			% \label{EMS-1840-JS-DAY-3LETTERMONTH-YEAR}
			
			\begin{description}
				\item[Print Sources]
			\end{description}

			\begin{tabular}{p{0.5in}p{1in}p{0.5in}p{1in}}
				JSP	 	& ---		& JSR 		& --- \\
				DC	 	& ---		& HC 		& --- \\
				HDDC 	& ---		& TCHC 		& --- \\
			\end{tabular}
			% HDDC = woodford's DC diss; JSR = Marquardt's JS Revelations; TCHC = Vogel HC
			
			\begin{description}
				\item[Online Access] \href{https://www.josephsmithpapers.org/paper-summary/memorial-to-nauvoo-high-council-18-june-1840/1}{Here}
				% \item[Analysis] \hyperref[XX]{Here}
			\end{description}
			
			\begin{description}
				\item[Background]
			\end{description}
			
				% It might also be useful to give a two sentence context of the period: e.g., "This speech was given in Nauvoo to a small congregation one month prior to the destruction of the Nauvoo Expositor. These and other tensions may have been on the prophet's mind when he gave the speech."
			
			\begin{description}
				\item[Original Source]
			\end{description}
			
				There is another version of this document which has important changes, especially in the part about translating the documents; see the notes in JSP. I should track this down and get the changes.
			
				% brief description of original source material, with reference to JSP or somewhere else for more info
				% Background material: it might be helpful to have a note that says whether a given document was presented as a speech, was written in a journal, etc.
				% Also a comment here about how reliable the source might be, or what shape it is in, if there are large omissions or parts that are hard to understand, and so on.
			
			\begin{description}
				\item[Text]
			\end{description}
			
					\hypertarget{EMS-1840-JS-18-JUN-1840-a}%
				To
					\marginnote{a}%
				the Honorable the High Council of the Church of Jesus Christ of Latter Day Saints

				The memorial of Joseph Smith Jr. respectfully represents.

				That after the Church of Jesus Christ had been inhumanly as well as unconstitutionally expelled from their homes which they had secured to themselves in the State of Missouri, \sout{and having} <they> found a resting place in the State of Illinois altho very much scattered and at considerable distances from each other.

				That after the escape of your Memorialist from his enemies, he, (under the direction of the Authorities of the Church) took such Steps as has secured to the church the present Locations viz the Town plat of Nauvoo and lands in the Iowa Territory

					\hypertarget{EMS-1840-JS-18-JUN-1840-b}%
				That
					\marginnote{b}%
				in order to secure said locations your Memorialist had to become responsible for the payments of the same and had to use considerable exertion in order to commence a settlement and a place of gathering for the Saints, but knowing that from the genius of the constitution of the church and for the well being of the saints that it was necessary so that the Constituted Authorities of the Church might assemble together to act as to legislate for the good of the whole society and that the saints might enjoy those priviledges which they could not by being scattered so wide apart, induced your Memorialist, to exert himself to the utmost, in order to bring about an object so necessary and so desireable to the Saints at large

					\hypertarget{EMS-1840-JS-18-JUN-1840-c}%
				That
					\marginnote{c}%
				under the then existing circumstances Your Memorialist had necessarily to engage in the temporalities of the Church--- which he has had to attend to the present time which has greatly engaged his mind and taken up much of his time.

				That your Memorialist feels it a duty which he owes to God as well as to the Church to give his attention more particularly to those things connected with the Spiritual welfare of the Saints, (which have now become a great people) so that they may be built up in their most holy faith and be eneabled to go on to perfection------

					\hypertarget{EMS-1840-JS-18-JUN-1840-d}%
				That
					\marginnote{d}%
				the church having erected an office where he can attend to the affairs of the church without distraction, he thinks and \sout{very} verily believes that \sout{in} the time has now come when he should devote himself exclusively to those things which relate to Spiritualities of the church and commence the work of translating the ejyptian Records--- the Bible--- and wait upon the Lord for Such revelations as may be suited to the condition and circumstances of the church and in order to attend to those things, prays that your honorable body will relieve him from the anxiety and trouble necessarily attendant on business transactions by appointing some \sout{one of the} Bishops to take charge of the City Plot and attend to the business transactions which have heretofore rested upon your Memorialists

					\hypertarget{EMS-1840-JS-18-JUN-1840-e}%
				That
					\marginnote{e}%
				should your Honors deem it, propper to do so, your memorialist would respectfully suggest, that he would have no means \sout{for} of support whatever and therefore would request that some one might be appointed to see that all his necessary wants be provided for as well as sufficient means or appropriations for a Clerk or Clerks which he may require to aid him in his important work

				Your Memorialist would further represent that as Elder H[enry] G. Sherwood is conversant with the affairs of the city plot, he thinks that he would be a suitable person to act as Clerk in that business and attend to the disposing of the remaining lots \&c \&c

					\hypertarget{EMS-1840-JS-18-JUN-1840-f}%
				Your
					\marginnote{f}%
				Memorialist would take this oppertunity of congratulating your Honorable body on the peace and Harmony which exists in the Church and for the good feelings which Seem to be manifest by all the Saints--- and hopes that inasmuch as every one devotes themselves for the good of the church and the spread of the Kingdom that the Choisest blessings of Heaven will be poured upon us, and that the Glory of the Lord will overshadow the inheritances of the Saints

				\begin{tightright}
					Joseph Smith Jr
				\end{tightright}

				Nauvoo June 18\supers{th.} 1840

				To the Hon

				The High Council of the Church of Jesus Christ Of \uline{Latter Day Saints}
				
		% -----
		\subsubsection{Minutes, 20 June 1840}
		% -----
			\label{EMS-1840-JS-20-JUN-1840}
			% \label{EMS-1840-JS-DAY-3LETTERMONTH-YEAR}
			
			\begin{description}
				\item[Print Sources]
			\end{description}

			\begin{tabular}{p{0.5in}p{1in}p{0.5in}p{1in}}
				JSP	 	& ---		& JSR 		& --- \\
				DC	 	& ---		& HC 		& --- \\
				HDDC 	& ---		& TCHC 		& --- \\
			\end{tabular}
			% HDDC = woodford's DC diss; JSR = Marquardt's JS Revelations; TCHC = Vogel HC
			
			\begin{description}
				\item[Online Access] \href{https://www.josephsmithpapers.org/paper-summary/minutes-20-june-1840/1}{Here}
				% \item[Analysis] \hyperref[XX]{Here}
			\end{description}
			
			\begin{description}
				\item[Background]
			\end{description}
			
				% It might also be useful to give a two sentence context of the period: e.g., "This speech was given in Nauvoo to a small congregation one month prior to the destruction of the Nauvoo Expositor. These and other tensions may have been on the prophet's mind when he gave the speech."
			
			\begin{description}
				\item[Original Source]
			\end{description}
			
				% brief description of original source material, with reference to JSP or somewhere else for more info
				% Background material: it might be helpful to have a note that says whether a given document was presented as a speech, was written in a journal, etc.
				% Also a comment here about how reliable the source might be, or what shape it is in, if there are large omissions or parts that are hard to understand, and so on.
			
			\begin{description}
				\item[Text]
			\end{description}
			
					\hypertarget{EMS-1840-JS-20-JUN-1840-a}%
				June
					\marginnote{a}%
				20\supers{th} 1840. The High Council of the Church of Jesus Christ of Latter Saints of Nauvoo Illinois met in council at the office [of] Joseph Smith jr.---

				1. A memmorial was presented to the council, by Elder R[obert] B. Thompson, from President Joseph Smith j\supersu{r} praying that the Council would relieve him from the temporalities of the Church which he necessarily had to engage in. In order that a location might be effected for the geathering of the saints.---

				That he felt it his duty to engage, more particularly, in the spiritual welfare of the Saints \& also, to the translating of the Egyptian Records--- the Bible--- \& wait upon the Lord for such revelations as may be suited to the condition \& circumstances of the Church.

					\hypertarget{EMS-1840-JS-20-JUN-1840-b}%
				That
					\marginnote{b}%
				in order to relieve him from the anxiety and troubles necessarily attendant on buisness--- transactions, they would appoint some one to take charge of the City Plot \& attend to the buisness transactions which have heretofore rested upon him.

				That when he is relieved from such temporal--- duties he would have no means of support whatever \& requested that some one might be appointed to see that all his necessary wants are provided for as well as sufficient means or appropriations for a clerk or clerks which <he> may require to aid him in his important work.

				The council relieved Presdent Joseph Smith jr according to \sout{this} the request of the memorial and appointed Elder H[enry] G. Sherwood to take charge of the City Plot and act as clerk in that buisness and also to attend to the disposing of the remaining lots and the buisness transactions which have heretofore rested upon him...%
					\footnote{%
						% \label{}%
						The JSP introduction to minutes at a 27 June 1840 conference: ``On 27 June 1840, Alanson Ripley, acting on behalf of JS, informed the Nauvoo high council that JS had vetoed actions the council took in its 20 June 1840 meeting, which JS apparently did not attend. At that meeting, the high council considered a memorial from JS requesting that he be relieved from his duties relating to temporal affairs in Nauvoo, Illinois, especially those involving land sales. Pursuant to JS's request, the high council appointed Henry G. Sherwood as a clerk for the sale of Nauvoo town lots and also directed Ripley to ensure that the temporal needs of JS and his counselors in the First Presidency were met. Although the high council's actions seem to meet the requests in JS's memorial, JS objected to these decisions. Minutes of a subsequent meeting suggest JS may have believed that the high council members had not sufficiently considered the heavy financial responsibility he bore from purchasing land in the Nauvoo area on behalf of the church or that they had not addressed pressing questions, such as how much Sherwood should be compensated for his service. JS also may have wanted the high council to specify where funding would come from to support him, his family, and the other members of the First Presidency. Whatever the case, Ripley relayed JS's veto, and the high council decided to rehear JS's memorial on \hyperref[EMS-1840-JS-3-JUL-1840]{3 July 1840}.''
						}
						
		% -----
		\subsubsection{Minutes, 3 July 1840}
		% -----
			\label{EMS-1840-JS-3-JUL-1840}
			% \label{EMS-1840-JS-DAY-3LETTERMONTH-YEAR}
			
			\begin{description}
				\item[Print Sources]
			\end{description}

			\begin{tabular}{p{0.5in}p{1in}p{0.5in}p{1in}}
				JSP	 	& ---		& JSR 		& --- \\
				DC	 	& ---		& HC 		& --- \\
				HDDC 	& ---		& TCHC 		& --- \\
			\end{tabular}
			% HDDC = woodford's DC diss; JSR = Marquardt's JS Revelations; TCHC = Vogel HC
			
			\begin{description}
				\item[Online Access] \href{https://www.josephsmithpapers.org/paper-summary/minutes-3-july-1840/1}{Here}
				% \item[Analysis] \hyperref[XX]{Here}
			\end{description}
			
			\begin{description}
				\item[Background]
			\end{description}
			
				% It might also be useful to give a two sentence context of the period: e.g., "This speech was given in Nauvoo to a small congregation one month prior to the destruction of the Nauvoo Expositor. These and other tensions may have been on the prophet's mind when he gave the speech."
			
			\begin{description}
				\item[Original Source]
			\end{description}
			
				% brief description of original source material, with reference to JSP or somewhere else for more info
				% Background material: it might be helpful to have a note that says whether a given document was presented as a speech, was written in a journal, etc.
				% Also a comment here about how reliable the source might be, or what shape it is in, if there are large omissions or parts that are hard to understand, and so on.
			
			\begin{description}
				\item[Text]
			\end{description}
			
					\hypertarget{EMS-1840-JS-3-JUL-1840-a}%
				July
					\marginnote{a}%
				3\supers{rd} \sout{1842} <1840>. The High Council of the Church of Jesus Christ of Latter Day Saints of Nauvoo Illinois met in Council at the office of Joseph Smith jr.

				1 The subject of the memmorial of Pres\supers{t} Joseph Smith j\supersu{r} was again taken up for a rehearing, according to the decision of the last Council (June 27) and the following resolutions were passed unanimously.

					\hypertarget{EMS-1840-JS-3-JUL-1840-b}%
				1\supers{st}
					\marginnote{b}%
				Resolved--- that we feel perfectly satisfied with the course taken by Joseph Smith j\supersu{r} \& feel a disposition as far as it is in our power to assist him so as to relieve him from the temporalities of the Church in order that he may devote his time more particularly to the spiritualities of the same, believeing by so doing we shall promote the good of the whole Church. But as he (Joseph Smith j\supers{r}) is held responsible for the payment of the City <Plot> and knowing no way to relieve him from that responsibility at present, we would request of him to act as treasurer for the City Plot and to whom those persons whom we may appoint to make sales of Lots and attend to the buisness affairs of the Church may at all times be responsible and make true and correct returns of all their proceedings as well as to account for all monies, properties \&c which may come into their hands.

					\hypertarget{EMS-1840-JS-3-JUL-1840-c}%
				Therefore
					\marginnote{c}%
				Resolved. That Elder Henry G Sherwood act as clerk for the same. That Bishop Alonson [Alanson] Ripley be appointed to provide for the wants of the Presidency and make such appropriations to them and to their Clerk or Clerks which they may require.

				Resolved. that the funds of the City Plot shall not be taken to provide for the Presidency or Clerks but that the Bishops be instructed to raise funds from other sources to meet the calls made on them.

				And mon[i]es recieved for lots shall be deposited in the hands of the treasurer to liquidate the debts of the City Plot.

					\hypertarget{EMS-1840-JS-3-JUL-1840-d}%
				Resolved
					\marginnote{d}%
				that the clerk shall have a stipulated sum for his services for transacting the buisness of the City Plot and that he recieve the sum of \uline{twenty}-\uline{six} dollars p\supersu{r} month for said services.---

				Resolved--- that the remaining items of buisness relative to the memmorial before the Council shall be laid over to be taken up again at the next Council.

				Adjourned

				\begin{tightright}
					Hosea Stout Clerk pro tem
				\end{tightright}
				
		% -----
		\subsubsection{Minutes, 11 July 1840}
		% -----
			\label{EMS-1840-JS-11-JUL-1840}
			% \label{EMS-1840-JS-DAY-3LETTERMONTH-YEAR}
			
			\begin{description}
				\item[Print Sources]
			\end{description}

			\begin{tabular}{p{0.5in}p{1in}p{0.5in}p{1in}}
				JSP	 	& ---		& JSR 		& --- \\
				DC	 	& ---		& HC 		& --- \\
				HDDC 	& ---		& TCHC 		& --- \\
			\end{tabular}
			% HDDC = woodford's DC diss; JSR = Marquardt's JS Revelations; TCHC = Vogel HC
			
			\begin{description}
				\item[Online Access] \href{https://www.josephsmithpapers.org/paper-summary/minutes-11-july-1840/1}{Here}
				% \item[Analysis] \hyperref[XX]{Here}
			\end{description}
			
			\begin{description}
				\item[Background]
			\end{description}
			
				% It might also be useful to give a two sentence context of the period: e.g., "This speech was given in Nauvoo to a small congregation one month prior to the destruction of the Nauvoo Expositor. These and other tensions may have been on the prophet's mind when he gave the speech."
			
			\begin{description}
				\item[Original Source]
			\end{description}
			
				% brief description of original source material, with reference to JSP or somewhere else for more info
				% Background material: it might be helpful to have a note that says whether a given document was presented as a speech, was written in a journal, etc.
				% Also a comment here about how reliable the source might be, or what shape it is in, if there are large omissions or parts that are hard to understand, and so on.
			
			\begin{description}
				\item[Text]
			\end{description}
			
					\hypertarget{EMS-1840-JS-11-JUL-1840-a}%
				July
					\marginnote{a}%
				11\supers{th} \sout{1842} 1840. High Council of the Church of Jesus Christ of Latter Day Saints of Nauvoo Illinois met in Council at the office of Joseph Smith jr.
				
				There was no buisness before the council President Joseph Smith j\supers{r} was present, who taught the Councellors things relative to their duty in Council. and give the \sout{folloing} following precedent for the Council to be guided by and ordered the same to be recorded to. wit---
				
					\hypertarget{EMS-1840-JS-11-JUL-1840-b}%	
				That
					\marginnote{b}%
				the Council should try no case without both parties being present or having had an opportunity to be present. neither should they hear one parties complaint before his case is brought up for trial--- neither should they suffer the character of any one to be exposed before the High Council without the person being present and ready to defend him or herself---that the minds of the Councellors be not prejudiced for or against any one whose case they may possibly have to act upon. Adjourned
				
				\begin{tightright}
					Hosea Stout Clk pro tem
				\end{tightright}
				
		% -----
		\subsubsection{Discourse, circa 19 July 1840}
		% -----
			\label{EMS-1840-JS-CIR-19-JUL-1840}
			% \label{EMS-1840-JS-DAY-3LETTERMONTH-YEAR}
			
			% --
			\paragraph{As Reported by Martha Jane Knowlton Coray--A}
			% --
				\label{EMS-1840-JS-CIR-19-JUL-1840-MC-A}
			
				\begin{description}
					\item[Print Sources]
				\end{description}

				\begin{tabular}{p{0.5in}p{1in}p{0.5in}p{1in}}
					JSP	 	& ---		& JSR 		& --- \\
					DC	 	& ---		& HC 		& --- \\
					HDDC 	& ---		& TCHC 		& --- \\
				\end{tabular}
				% HDDC = woodford's DC diss; JSR = Marquardt's JS Revelations; TCHC = Vogel HC
				
				\begin{description}
					\item[Online Access] \href{https://www.josephsmithpapers.org/paper-summary/discourse-circa-19-july-1840-as-reported-by-martha-jane-knowlton-coray-a/1}{Here}
					% \item[Analysis] \hyperref[XX]{Here}
				\end{description}
				
				\begin{description}
					\item[Background]
				\end{description}
				
					% It might also be useful to give a two sentence context of the period: e.g., "This speech was given in Nauvoo to a small congregation one month prior to the destruction of the Nauvoo Expositor. These and other tensions may have been on the prophet's mind when he gave the speech."
				
				\begin{description}
					\item[Original Source]
				\end{description}
				
					% brief description of original source material, with reference to JSP or somewhere else for more info
					% Background material: it might be helpful to have a note that says whether a given document was presented as a speech, was written in a journal, etc.
					% Also a comment here about how reliable the source might be, or what shape it is in, if there are large omissions or parts that are hard to understand, and so on.
				
				\begin{description}
					\item[Text]
				\end{description}
				
						\hypertarget{EMS-1840-JS-CIR-19-JUL-1840-MC-A-a}%
					it
						\marginnote{a}%
					is wisdom in God that the saints as fast as \sout{for} conveinient should come up and we should enter into as compact a city as possible. we are not grea[t]ly distressed nor ever will be--- a stake appointed--- this is the principle place of gathering The[r]efore as many as can come let the Brethren begin to roll in like clouds we will sell you lots and if you cant pay for them you shall [have] them whosoever will come let them come and partake of the poverty of Nauvoo freely and those \sout{who} who partake of our poverty shall [be?] greatly blessed yea exceedingly for if they come in times of prosperity will put on their long faced pharisaical visages and blow us up
						\hypertarget{EMS-1840-JS-CIR-19-JUL-1840-MC-A-b}%
					and
						\marginnote{b}%
					this proclamation goes forth from this hour bring every thing you can bring and build the house of God and we will have a tremenduous City which shall reverberate afar--- afar--- and we will have a great observatory great great high wacth tower then comes all the ancient records dig them up thus shall the poor be fed by the curious who where the Saints g[ather] is Zion which <the> ene [rig]hteous will build u[p] a place of safety for the[i]r [c]hildren and those who [do] not come in due season shall scarce escape wi[t]h <their lives> those that cease to <be> saviours of men they shall be troden under the feet of their enemies for their transgression instance Reed peck
						\hypertarget{EMS-1840-JS-CIR-19-JUL-1840-MC-A-c}%
					The
						\marginnote{c}%
					redemption of Zion is the salvation of \sout{Zion} this country which is north and south America some consider the olive trees to be Jackson Co but one is not 12 \sout{six} some have a temple begun 12 olive T<rees> are 12 Stakes [\textit{illegible}] the seed of 12 olive [tree]s are scattered it [\textit{illegible}] wake up the na[tio]ns--- The nations will [wa]r with each other [w]hile the Saints are [b]uilding the 12 I had [a] vision on the other side of the river that my armies were about ne [me] and I ran from brethren--- who were shooting at me swan [swam] the river and found safety on this side Go tell all my servants who are the Strength of my house \&c this time this \sout{wil} nation will be trembling then shall foreign Saints come and all the saints shall come [an]d fight for God.
						\hypertarget{EMS-1840-JS-CIR-19-JUL-1840-MC-A-d}%
					they
						\marginnote{d}%
					\phantom{ }[s]hall come to the land [of] my vinyard [T]hen will Kings queens [r]ulers come and help buil[d] up Zion reason[in]g cant [he]lp themselves------ therefore go to with your might build build and you shall inherit it long and be blessed and the rising generation shall awake the nations make k[i]ngdoms tremble and the earth shake to the centre and God will bless the great rulers and cast a vot[e] against van buren \sout{before} Jerusalem must be built up Zion also [b]efore the Coming of Christ this will take nea[r] 60 years. all who can coolly and deliberately dispose of their property and come and give \uline{their} property to the poor an[d] worship God O bretheren come---
						\hypertarget{EMS-1840-JS-CIR-19-JUL-1840-MC-A-e}%
					I
						\marginnote{e}%
					prophcy [prophesy] this state shall become greate and <a> mighty mountain--- and this will be the greatest city in the world. curse the man who says you are a mean man because you do not believe as I do--- all libe[r]al men o come up [t]o nauvoo and help us to build up the City of God--- \sout{Go} the oposition from without does not compare with that within Moses led the children of Israel. the calamatiyes which happened to them \sout{is} was in consequence of rebeling against moses--- rbelious [rebellious] here also--- when I disbelheve in Joseph I will ask God for another is God a monarch if you disbelieve an honorable way is to say I would rather be damned than to follow self exaltation first cause of sin--- prayer--- first <thy kingdom come--- thy will> be done on earth as in \sout{the} \uline{H}[\uline{e}]\uline{aven}--- what would we do singing forever an[d] ever ans.--- starve to death
						\hypertarget{EMS-1840-JS-CIR-19-JUL-1840-MC-A-f}%
					the
						\marginnote{f}%
					Kingdom--- O[f]--- G[od]--- is like purified Elements--- like \sout{2} 1000ds and 1000ds Of worlds have been to effect the purif[y]ing is to pray for the Kingdom come \&c forgive as we are forgiven God had no hand in put[t]ing men in prison for debt but devils incarnate--- were the flame to roll as fast as we want it to it would crush half the church before it touched the world--- righteous principles--- instanse--- a man that has the Kingdom and its righteousness knows how to manage property what the bishops do is by d[i]rection of the presidency sectarians worship the 12 apostles---
						\hypertarget{EMS-1840-JS-CIR-19-JUL-1840-MC-A-g}%
					I
						\marginnote{g}%
					don't think that a man who says he will [\textit{illegible}] shall come forom [from] all parts of the world and I prophecy that pleasure parties come from england to see ther mamoth and like Sheba--- and schools houses and great men's sons shall board here and crave the privelidges of educating there children here and then I prophecy shall the Saints chink the money in their pockets, which they shall give
						\hypertarget{EMS-1840-JS-CIR-19-JUL-1840-MC-A-h}%
					And
						\marginnote{h}%
					obligate myself to build as great a house as ever Solomen did \sout{and} if the church will back me up--- and it shall impov[e]rish no one but enrich thousands The time comes when these Saints shall ride proudly over the mountains of Mis[so]uri and no gentile \sout{dog} dog shall dare lift a tongue but shall l[i]ck up the dust of their feet in the name of Jesus May many here realised this and see it with their eyes I myself also is my prayer--- Amen---
					
			% --
			\paragraph{As Reported by Martha Jane Knowlton Coray--B}
			% --
				\label{EMS-1840-JS-CIR-19-JUL-1840-MC-B}
			
				\begin{description}
					\item[Print Sources]
				\end{description}

				\begin{tabular}{p{0.5in}p{1in}p{0.5in}p{1in}}
					JSP	 	& ---		& JSR 		& --- \\
					DC	 	& ---		& HC 		& --- \\
					HDDC 	& ---		& TCHC 		& --- \\
				\end{tabular}
				% HDDC = woodford's DC diss; JSR = Marquardt's JS Revelations; TCHC = Vogel HC
				
				\begin{description}
					\item[Online Access] \href{https://www.josephsmithpapers.org/paper-summary/discourse-circa-19-july-1840-as-reported-by-martha-jane-knowlton-coray-b/1}{Here}
					% \item[Analysis] \hyperref[XX]{Here}
				\end{description}
				
				\begin{description}
					\item[Background]
				\end{description}
				
					% It might also be useful to give a two sentence context of the period: e.g., "This speech was given in Nauvoo to a small congregation one month prior to the destruction of the Nauvoo Expositor. These and other tensions may have been on the prophet's mind when he gave the speech."
				
				\begin{description}
					\item[Original Source]
				\end{description}
				
					% brief description of original source material, with reference to JSP or somewhere else for more info
					% Background material: it might be helpful to have a note that says whether a given document was presented as a speech, was written in a journal, etc.
					% Also a comment here about how reliable the source might be, or what shape it is in, if there are large omissions or parts that are hard to understand, and so on.
				
				\begin{description}
					\item[Text]
				\end{description}
				
						\hypertarget{EMS-1840-JS-CIR-19-JUL-1840-MC-B-a}%
					A
						\marginnote{a}%
					few Item from a discourse delivered by Joseph Smith July 19--- 1840

					Read a chap in [Ezekiel] ed concluding with this saying and when all these things come to pass and Lo they will come then shall you know that a Prophet hath been among you

					Afterwards read the parable of the 12 olive trees and said speaking of the Land of Zion \sout{that} It \sout{consisted} consists of all N. \& S America but that any place where the Saints gather is Zion which every righteous man will build up for a place of safety for his children \sout{that} The olive trees are 12 stakes which are yet to be built not the Temple in Jackson as some suppose for while the 12 \sout{olive} stakes are being built we will be at peace but the Nations of the Earth will be at war:

						\hypertarget{EMS-1840-JS-CIR-19-JUL-1840-MC-B-b}%
					our
						\marginnote{b}%
					cry from the 1\supersu{st} has been for peace and we will continue pleading like the Widow at the feet of the unjust judge but we may plead at the feet of Majistrates and at the feet of Judges At the feet of Governors and at the feet of senotors \& at the feet of \sout{the} Pre[s]idents for 8 years it will be of no avail we shall find no favor in any <of the> courts of this government

						\hypertarget{EMS-1840-JS-CIR-19-JUL-1840-MC-B-c}%
					The
						\marginnote{c}%
					redemption of Zion is the redemption of all N \& S America and those 12 stake must be built up before the redemption of Zion can take place and those who refuse to gather and build when they are comanded to do so cease to be Saviours of men and are thenceforth good for nothing but shall be cast out and trodden underfeet of men for their transgression as Reed Peck was when he aplied in the name of an apostate for buisness in a store in Quincy They told him that they wanted no apostates round them and showed him the door At this same store the Authorities of this Churc<h> could have obtained almost any amount of Credit they could have asked---

						\hypertarget{EMS-1840-JS-CIR-19-JUL-1840-MC-B-d}%
					*We
						\marginnote{d}%
					shall build the Zion of the Lord in peace untill the servants of that Lord shall begin to lay the found[a]tion of a great and high watch Tower and then shall they begin to say within themselves what need hath my Lord of this tower seeing this is a time of peace \&c---

						\hypertarget{EMS-1840-JS-CIR-19-JUL-1840-MC-B-e}%
					Then
						\marginnote{e}%
					the Enemy shall come as a thief in the night and [\textit{one blank line}] scatter the servants abroad when the seed of these 12 Olive trees are scattered abroad they will wake up the Nations of the whole Earth Eeven this Nation will be on the very verge of crumbling to peices and tumbling to the ground and when the constitution is upon the brink of ruin this people will be the Staff up[on] which the Nation shall lean and they shall bear \sout{away} the constitution away from the <very> verge of destruction--- Then shall the Lord say Go tell all my servants who are the strength of mine house my young men and middle aged \&c Come to the Land of my vineyard and fight the battle of the Lord---
						\hypertarget{EMS-1840-JS-CIR-19-JUL-1840-MC-B-f}%
					Then
						\marginnote{f}%
					the Kings \& Queens shall come then the rulers of the Earth shall come then shall all saints come yea the Foreign saints shall come to fight for the Land of my vineyard for in this thing Shall be their safety and they will have no power to choose but will come as a man fleeeth from a sudden des destruction--- But before this the time shall be \sout{when} these who are now my friends shall become my enemies and shall seek to take my life and \sout{shall be m} there are those now before me who will more furiously pursue me \sout{and} more dilligently seek \sout{to} my life and be more b[l]ood thirsty upon my track \sout{thus} <than> ever were the Misouri\sout{an} Mobbers you say among yourselves as did them of old time is it I \& is it I but I know these things by the visions of the Almighty---

						\hypertarget{EMS-1840-JS-CIR-19-JUL-1840-MC-B-g}%
					But
						\marginnote{g}%
					brethren come ye yea come all of you who can come and go to with your mights and build up the cities of the Lord and whosoever will let him come and partake of the poverty of Nauvoo freely for those who partake of her poverty shall also \sout{shall also} partake of her prosperity and it is now wisdom in God that we should enter into as compact a city as posible For Zion and Jerusalem \sout{and} must both be built up before the coming of Christ \sout{This} \sout{will} \sout{near} \sout{a} \sout{half} \sout{of} \sout{a} \sout{century} How long will it take to do this 10 years yes more than 40 years will pass before this work will be accomplished and when these cities are built then shall the coming of the son of man be

						\hypertarget{EMS-1840-JS-CIR-19-JUL-1840-MC-B-h}%
					Now
						\marginnote{h}%
					let all who can coolly and deliberately dispose of their property come up and give of their substance to the [poor?] that the hearts of the poor may be comforted and all may worship god together in holiness of heart come brethren come all of you--- And I prophecy in the name of the Lord that the state of Illinois shall become a great \sout{mountain} and mighty mountain as [a] city set upon a hill that cannot be hid and a great [candle?] that giveth light to the world \sout{and} The city of Nauvoo als[o] shall become the Greatest city in the whole world---

						\hypertarget{EMS-1840-JS-CIR-19-JUL-1840-MC-B-i}%
					curse
						\marginnote{i}%
					that man who says to his neighbor you are a mean man because you do not believe as I do I now invite all liberall minded men to come up to Nauvoo and help to build up the city of our God We are not greatly distressed \sout{No} no nor ever will be This is the principle place of gathering therefore let the brethren begin to roll in like clouds and we will sell you lots if you are able to pay for them and \sout{of} if not you shall have them without money and without price

						\hypertarget{EMS-1840-JS-CIR-19-JUL-1840-MC-B-j}%
					The
						\marginnote{j}%
					greater blessing is unto those who come in times of adversity. For many will come to us in times of prosperity that will stand at the corners of the streets saying with long pharisaical faces to those that come after thems dont go near Bro Joseph dont go near the authorities of the church for they will pick your pockets they will rob you of all your money Thus will they breed in our midst a spirit of dissatisfaction and distrust that will end in persecution and distress------

						\hypertarget{EMS-1840-JS-CIR-19-JUL-1840-MC-B-k}%
					Now
						\marginnote{k}%
					from this hour bring every thing you can bring and build a Temble [Temple] unto the Lord a house into the mighty God of Jacob. We will build upon the top of this Temple a great observatory a great and high watch tower and in the top thereof we will Suspend a tremendous bell that when it is rung shall rouse the inhabitants of [Fort] Madison \sout{and} wake up the people of Warsaw and sound in the ears of men [in] carthage Then comes the ancient records yea all of them dig them yes bring them forth speedily

						\hypertarget{EMS-1840-JS-CIR-19-JUL-1840-MC-B-l}%
					Then
						\marginnote{l}%
					shall the poor be fed by the curious who shall come from all parts of the world to see this wonderful temple Yea I I prophecy that pleasure parties shall come from England to see the Mamoth and like the Queen of Sheba shall say the half never was told them. School houses shall be built here and High schools shall be established and the great men of the [earth] shall send their sons here to board while they are receiving their education among us and even Noblemen shall crave the priviledge of educating their children with us and these poor saints shall chink in their pockets the money of these proud men received from such as come and dwell with us

						\hypertarget{EMS-1840-JS-CIR-19-JUL-1840-MC-B-m}%
					Now
						\marginnote{m}%
					brethren I obligate myself to build as great a temple as ever solomon did if the church will back me up. moreover it shall not impoverish any man but enrich thousands and I prophecy that the time shall be when these saints shall ride proudly over the mountains of Missouri and no Gentile dog nor Missouri dog Shall dare lift a tongue against them but will lick up the dust from beneath their feet and I pray the father that many here may realize this and see it with their eyes \sout{and} And if it should be (Stretching his hand towards the place and in a melancholly tone that made all hearts tremble) will of God that I might live to behold that temple completed and finished from the foundation to the top stone I will say Oh Lord it is enough Lord let thy servant depart in peace which is my ernest prayer in the name of the L Jesus Amen
					
		% -----
		\subsubsection{Letter to William W. Phelps, 22 July 1840}
		% -----
			\label{EMS-1840-JS-22-JUL-1840}
			% \label{EMS-1840-JS-DAY-3LETTERMONTH-YEAR}
			
			\begin{description}
				\item[Print Sources]
			\end{description}

			\begin{tabular}{p{0.5in}p{1in}p{0.5in}p{1in}}
				JSP	 	& ---		& JSR 		& --- \\
				DC	 	& ---		& HC 		& --- \\
				HDDC 	& ---		& TCHC 		& --- \\
			\end{tabular}
			% HDDC = woodford's DC diss; JSR = Marquardt's JS Revelations; TCHC = Vogel HC
			
			\begin{description}
				\item[Online Access] \href{https://www.josephsmithpapers.org/paper-summary/letter-to-william-w-phelps-22-july-1840/1}{Here}
				% \item[Analysis] \hyperref[XX]{Here}
			\end{description}
			
			\begin{description}
				\item[Background]
			\end{description}
			
				% It might also be useful to give a two sentence context of the period: e.g., "This speech was given in Nauvoo to a small congregation one month prior to the destruction of the Nauvoo Expositor. These and other tensions may have been on the prophet's mind when he gave the speech."
			
			\begin{description}
				\item[Original Source]
			\end{description}
			
				% brief description of original source material, with reference to JSP or somewhere else for more info
				% Background material: it might be helpful to have a note that says whether a given document was presented as a speech, was written in a journal, etc.
				% Also a comment here about how reliable the source might be, or what shape it is in, if there are large omissions or parts that are hard to understand, and so on.
			
			\begin{description}
				\item[Text]
			\end{description}
			
				\begin{tightright}
						\hypertarget{EMS-1840-JS-22-JUL-1840-a}%
					Nauvoo
						\marginnote{a}%
					Hancock Co Ills
					
					July 22\supers{nd.} 1840
				\end{tightright}

				Dear Brother [William W.] Phelps

				I must say that it is with no ordinary feelings I endeavour to write a few lines to you in answer to yours of the 29\supers{th.} Ultimo, at the same time I am rejoiced at the priveledge granted me. You may in some measure realise what my feelings, as well as Elder [Sidney] Rigdon's \& Bro Hyrum [Smith]'s were when we read your letter, truly our hearts were melted into tenderness and compassion when we assertained your resolves \&c

					\hypertarget{EMS-1840-JS-22-JUL-1840-b}%
				I
					\marginnote{b}%
				can assure you I feel a disposition to act on your case in a manner that will meet the approbation of Jehovah (whose servant I am) and agreeably to the principles of truth and righteousness which have been revealed and inasmuch as long-suffering patience and mercy have ever characterized the dealings of our heavenly Father towards the humble and penitent, I feel disposed, to copy the example and cherish the same principles, by so doing be a savior of my fellow men

					\hypertarget{EMS-1840-JS-22-JUL-1840-c}%
				It
					\marginnote{c}%
				is true, that we have suffered much in consequence of your behavior---\uline{the} \uline{cup} \uline{of} \uline{gall} \uline{already} \uline{full} \uline{enough} for mortals to drink, was indeed \uline{filled} to \uline{overflowing} when you turned against us: One with whom we had oft taken sweet council together, and enjoyed many refreshing seasons from the Lord ``Had it been an enemy we could have borne it'' [``]In the day that thou stoodest on the other side, in the day when Strangers carried away captive his forces, and foreigners entered into his gates and cast lots upon Far West even thou wast as one of them. But thou shouldst not have looked on the day of thy brother, in the day that he became a stranger neither shouldst thou have spoken proudly in the day of distress''

					\hypertarget{EMS-1840-JS-22-JUL-1840-d}%
				However
					\marginnote{d}%
				the cup has been drunk, the will of our heavenly Father has been done, and we are yet alive for which we thank the Lord. And having been delivered from the hands of wicked men by the mercy of our God, we say it is your privilidge to be delivered from the power of the Adversary--- be brought into the liberty of God's dear children, and again take your stand among the saints of the Most High, and by diligence humility and love unfeigned, commend yourself to our God and your God and to the church of Jesus Christ

					\hypertarget{EMS-1840-JS-22-JUL-1840-e}%
				Believing
					\marginnote{e}%
				your confession to be real and your repentance genuine, I shall be happy once again to give you the right hand of fellowship, and rejoice over the returning prodigal.

				Your letter was read to the saints last sunday and an expression of their feeling was taken, when it was unanimously resolved that W. W. Phelps should be received into fellowship.

				\begin{tightcenter}
					``Come on dear Brother since the war is past,

					For friends at first are friends again at last.''
				\end{tightcenter}

				\begin{tightright}
					Yours as Ever

					Joseph Smith Jr
				\end{tightright}
				
		% -----
		\subsubsection{Letter to Oliver Granger, between circa 22 and circa 28 July 1840}
		% -----
			\label{EMS-1840-JS-CIR-22-CIR-28-JUL-1840}
			% \label{EMS-1840-JS-DAY-3LETTERMONTH-YEAR}
			
			\begin{description}
				\item[Print Sources]
			\end{description}

			\begin{tabular}{p{0.5in}p{1in}p{0.5in}p{1in}}
				JSP	 	& ---		& JSR 		& --- \\
				DC	 	& ---		& HC 		& --- \\
				HDDC 	& ---		& TCHC 		& --- \\
			\end{tabular}
			% HDDC = woodford's DC diss; JSR = Marquardt's JS Revelations; TCHC = Vogel HC
			
			\begin{description}
				\item[Online Access] \href{https://www.josephsmithpapers.org/paper-summary/letter-to-oliver-granger-between-circa-22-and-circa-28-july-1840/1}{Here}
				% \item[Analysis] \hyperref[XX]{Here}
			\end{description}
			
			\begin{description}
				\item[Background]
			\end{description}
			
				% It might also be useful to give a two sentence context of the period: e.g., "This speech was given in Nauvoo to a small congregation one month prior to the destruction of the Nauvoo Expositor. These and other tensions may have been on the prophet's mind when he gave the speech."
			
			\begin{description}
				\item[Original Source]
			\end{description}
			
				% brief description of original source material, with reference to JSP or somewhere else for more info
				% Background material: it might be helpful to have a note that says whether a given document was presented as a speech, was written in a journal, etc.
				% Also a comment here about how reliable the source might be, or what shape it is in, if there are large omissions or parts that are hard to understand, and so on.
			
			\begin{description}
				\item[Text]
			\end{description}
			
					\hypertarget{EMS-1840-JS-CIR-22-CIR-28-JUL-1840-a}%
				Brother
					\marginnote{a}%
				\phantom{ }[Oliver] Granger % like elsehwere, this \phantom{ } is because of what needs to come after the marginal note

				D\supers{r} Sir

				It was with great pleasure I received your's and Bro Richards [Levi Richards's] Letter dated New york June 23\supersu{rd}\supers{.} 1840 and was very happy to be informed of your safe arrival in that place and your probability of success, and I do hope that your anticipations will be realized and that you will be able to free the Lords House from all incumbrances, and be prospered in all your undertakings for the benefit of the Church, and pray that while you are exerting your influence to bring about an object so desireable, that the choicest blessings of heaven may rest down upon you.
					\hypertarget{EMS-1840-JS-CIR-22-CIR-28-JUL-1840-b}%
				While
					\marginnote{b}%
				you are endeavoring to do so and attending to the duties laid upon you by the Authorities of the Church in this place, I am sorry to be informed not only in your letter but from other respectable sources of the strange conduct pursued in Kirtland by Elder Alman Babbit [Almon Babbitt]; I am indeed Surprised that a man having the experience which Bro Babbit has had should take any steps whatever calculated to destroy the confidence of the brethren in the presidency or any of the Authorities of the church. In order to conduct the affairs of the kingdom in righteousness it is all important, that the most perfect harmony kind feeling, good understanding and confidence should exist in the hearts of all the brethren. and that true Charity---love one towards another, should characterize all their proceedings.
					\hypertarget{EMS-1840-JS-CIR-22-CIR-28-JUL-1840-c}%
				If
					\marginnote{c}%
				there are any uncharitable feelings, any lack of confidence, then pride and arrogancy and envy will soon be manifested and confusion must inevitably prevail and the Authorities of the Church set at nought; and under such circumstances Kirtland cannot rise and free herself from <the> captivity in which she is held and become a place of safety for the saints nor can the blessings of Jehovah rest upon her. If the saints in Kirtland deem me unworthy of their prayers when they assemble together, and neglect to bear me up at a throne of heavenly grace, it is a strong and convincing proof to me that they have not the spirit of God.

					\hypertarget{EMS-1840-JS-CIR-22-CIR-28-JUL-1840-d}%
				If
					\marginnote{d}%
				the revelations we have received are true, who is to lead the people? If the keys of the kingdom have been committed to my hands, who shall open out the mysteries thereof. As long as my brethren stand by me and encourage me I can combat the predjudices of the world and can bear the contumely and abuse \sout{of the world} with joy but when my brethren stand aloof--- when they begin to faint and endeavour to retard my progress and enterprise then I feel to mourn but am no less determined to prosecute my task, being confident that altho my earthly friends may fail and even turn against me, yet my heavenly father will bear me off triumphant.
					\hypertarget{EMS-1840-JS-CIR-22-CIR-28-JUL-1840-e}%
				However
					\marginnote{e}%
				I hope that even in Kirtland, their are some who do not make a man an [``]offender for a word'' but are disposed to stand forth in defence of righteousness and truth and attend to every duty enjoined upon them and who will have wisdom to direct them against any movement or influence calculated to bring confusion and disorder into the camp of Israel, and to discern between the spirit of truth and the spirit of error.

					\hypertarget{EMS-1840-JS-CIR-22-CIR-28-JUL-1840-f}%
				It
					\marginnote{f}%
				would be gratifying to my mind to see the saints in Kirtland flourish, but think the time has not yet come and I assure you it never will until a different order of things be established and a different spirit be manifested. When confidence is restored, when pride shall fall and every aspiring mind be clothed with humility as with a garment and selfishness give place to benevolence and charity, and a united determination to live by every word which proceedeth out of the mouth of the Lord is observable, then and not till then can peace and order, and love prevail

					\hypertarget{EMS-1840-JS-CIR-22-CIR-28-JUL-1840-g}%
				It
					\marginnote{g}%
				is in consequence of aspiring men that Kirtland has been forsaken. How frequently has your humble servant been envied in his office by such characters who endeavoured to raise themselves to power at my expense, and seeing it impossible to do so, resorted to foul slander and abuse and other means to effect my overthrow; such characters have ever been the first to cry out against the presidency, and publish their faults and foibles to four winds of heaven.

					\hypertarget{EMS-1840-JS-CIR-22-CIR-28-JUL-1840-h}%
				I
					\marginnote{h}%
				cannot forget the treatment I received in the house of my friends, these things continually roll across my mind and cause me much sorrow of heart, and when I think that others who have lately come into the church should be led to Kirtland instead of to this place by Elder Babbit, and having their confidence in the Authorities lessened by such observations as he (Elder Babbit) has thought propper to make, as well as hearing all the false reports and exaggerated accounts of our enemies, I must say that I feel grieved in spirit, and cannot tolerate such proceedings neither will I, but will endeavour to disabuse the minds of the saints and break down all such unhallowed proceedings.

					\hypertarget{EMS-1840-JS-CIR-22-CIR-28-JUL-1840-i}%
				It
					\marginnote{i}%
				was something new to me when I heard there had been secret meetings held in the Lords house, and that some of my friends---faithful brethren, men enjoying the confidence of the church should be locked out. Such like proceedings are not calculated to promote union or peace but to engender strife and will be a curse instead of a blessing: To those who are young in the work I know they are calculated to and must be injurious to them. Those who have had experience and who should know better, than to reflect on their brethren, there is no excuse for them. If Bro Babbit and the other brethren wish to reform the Church and come out and make a stand against sin \& speculation \&c \&c; they must use other weapons than lies, or their object can never be effected, and their labors will be given to the house of the stranger rather than to the house of the Lord

					\hypertarget{EMS-1840-JS-CIR-22-CIR-28-JUL-1840-j}%
				The
					\marginnote{j}%
				proceedings of Bro Babbit were taken into consideration at a meeting of the Church at this place, when it was unanimously resolved that fellowship should be withdrawn from him until he make satisfaction for the conduct he has pursued of which circumstance I wish you to apprize him of without delay and demand his license.

				D\supers{r} Sir I wish you to stand in your lot and keep the station which was given you by revelation and the authorities of the Church; attend to the affairs of the Church with diligence and then rest assured on the blessings of heaven: It is binding on you to act as president of the Church in Kirtland until you are removed by the same Authority which put you in, and I do hope, their will be no cause for opposition. but that good feelin[g]s will be manifested in future by all the brethren

					\hypertarget{EMS-1840-JS-CIR-22-CIR-28-JUL-1840-k}%
				Bro
					\marginnote{k}%
				Burdicks [Thomas Burdick's] letter to Bro Hyrum [Smith] was duly received for which he has our best thanks, It was indeed an admirable letter and worthy of its author the sentiments express'd were in accordance with the spirit of the gospel and the principles correct. I am glad that Bro Richards has continued with you and hope he has been of some service to you,--- give my love to him

				Our prospects in this place continue good. considerable numbers have come in this spring.--- There were some bickerings respecting your conduct soon after your departure but they have all blown over, and I hope there will never be any occasion for any more, but that you will commend yourself to God and to the saints by a virtuous walk and holy conversation

				I had a letter from W[illiam] W. Phelps a few days ago informing me of his desire to come back to the Church if we would accept of him, he appears very humble and is willing to make every satisfaction that Saints or God may require.

				We expect to have an edition of the book of Mormon printed by the first of September it \sout{has} is now being sterotyped in Cincinnatti

				\begin{tightright}
					I rem[ain] \&c \&c

					Joseph Smith Jr
				\end{tightright}
				
		% -----
		\subsubsection{Discourse, 30 July 1840, as Reported by John Smith}
		% -----
			\label{EMS-1840-JS-30-JUL-1840}
			% \label{EMS-1840-JS-DAY-3LETTERMONTH-YEAR}
			
			\begin{description}
				\item[Print Sources]
			\end{description}

			\begin{tabular}{p{0.5in}p{1in}p{0.5in}p{1in}}
				JSP	 	& ---		& JSR 		& --- \\
				DC	 	& ---		& HC 		& --- \\
				HDDC 	& ---		& TCHC 		& --- \\
			\end{tabular}
			% HDDC = woodford's DC diss; JSR = Marquardt's JS Revelations; TCHC = Vogel HC
			
			\begin{description}
				\item[Online Access] \href{https://www.josephsmithpapers.org/paper-summary/discourse-30-july-1840-as-reported-by-john-smith/1}{Here}
				% \item[Analysis] \hyperref[XX]{Here}
			\end{description}
			
			\begin{description}
				\item[Background]
			\end{description}
			
				% It might also be useful to give a two sentence context of the period: e.g., "This speech was given in Nauvoo to a small congregation one month prior to the destruction of the Nauvoo Expositor. These and other tensions may have been on the prophet's mind when he gave the speech."
			
			\begin{description}
				\item[Original Source]
			\end{description}
			
				% brief description of original source material, with reference to JSP or somewhere else for more info
				% Background material: it might be helpful to have a note that says whether a given document was presented as a speech, was written in a journal, etc.
				% Also a comment here about how reliable the source might be, or what shape it is in, if there are large omissions or parts that are hard to understand, and so on.
			
			\begin{description}
				\item[Text]
			\end{description}
			
				Joseph said that inasmuch as the Saints will leave of speaking Evil of one another and not speake Evil of the Seer which <is> Speaking Evil of the ar[c]hangel or the holy Kees which he held and the Archangel and if the church would ceace from these Bickerings and murmuring and be of one mind the Lord would Visit\sout{ing} them with health and Every needed good for said he this is the Voice of the Spirit. furthermore <if> the Saints are sick or have Sickness in their Families and the Elders do not prevail every family should [get] Power by fasting and prayer and annointing with oil and continue So to Do their \sout{their} sick shall all be healed this also is the voice of the Spirit [p. [3]]
				
		% -----
		\subsubsection{Letter to John C. Bennett, 8 August 1840}
		% -----
			\label{EMS-1840-JS-8-AUG-1840}
			% \label{EMS-1840-JS-DAY-3LETTERMONTH-YEAR}
			
			\begin{description}
				\item[Print Sources]
			\end{description}

			\begin{tabular}{p{0.5in}p{1in}p{0.5in}p{1in}}
				JSP	 	& ---		& JSR 		& --- \\
				DC	 	& ---		& HC 		& --- \\
				HDDC 	& ---		& TCHC 		& --- \\
			\end{tabular}
			% HDDC = woodford's DC diss; JSR = Marquardt's JS Revelations; TCHC = Vogel HC
			
			\begin{description}
				\item[Online Access] \href{https://www.josephsmithpapers.org/paper-summary/letter-to-john-c-bennett-8-august-1840/1}{Here}
				% \item[Analysis] \hyperref[XX]{Here}
			\end{description}
			
			\begin{description}
				\item[Background]
			\end{description}
			
				% It might also be useful to give a two sentence context of the period: e.g., "This speech was given in Nauvoo to a small congregation one month prior to the destruction of the Nauvoo Expositor. These and other tensions may have been on the prophet's mind when he gave the speech."
			
			\begin{description}
				\item[Original Source]
			\end{description}
			
				% brief description of original source material, with reference to JSP or somewhere else for more info
				% Background material: it might be helpful to have a note that says whether a given document was presented as a speech, was written in a journal, etc.
				% Also a comment here about how reliable the source might be, or what shape it is in, if there are large omissions or parts that are hard to understand, and so on.
			
			\begin{description}
				\item[Text]
			\end{description}
			
				\begin{tightright}
						\hypertarget{EMS-1840-JS-8-AUG-1840-a}%
					Nauvoo
						\marginnote{a}%
					Hancock Co Ill

					Aug 8\supers{th.} 1840
				\end{tightright}

				Dear Sir

				Yours of the 25\supers{th.} Ultimo addressed to Elder [Sidney] Rigdon \& myself is received for which you have our thanks and to which I shall feel great pleasure in replying. Although I have not the pleasure of your acquaintance, yet from the kindness manifested towards our people when in bondage and oppression, and from \sout{the} the frank and noble mindedness br[e]athed in your letter, I am brot, to the conclusion that you are a friend to suffering humanity \& Truth.

				To those who have suffered so much abuse and borne the cruelties and insults of wicked men so long on account of those principles which we have been instructed to teach to the world a feeling of sympathy and kindness is something like the refreshing breese and cooling stream at the present season of the year and are I assure you duly appreciated by us

					\hypertarget{EMS-1840-JS-8-AUG-1840-b}%
				It
					\marginnote{b}%
				would afford me much pleasure to see you at this \sout{time} place, and from the desire you express in your letter to move to this place I hope I shall soon have that satisfaction.

				I have no doubt but you would be of great service to this community in practicing your profession as well as those other abilities of which you are in possession. Since to devote your time and abilities in the cause of truth and a suffering people may not be the means of exalting you in the eyes of this generation or securing you the riches of the world yet, by so doing you may rely on the approval of Jehovah ``That blessing which maketh rich and addeth not sorrow''

					\hypertarget{EMS-1840-JS-8-AUG-1840-c}%
				Through
					\marginnote{c}%
				the tender mercies of our God we have escaped the hands of those who sought our overthrow and have secured locations in this state and in the Territory of Iowa. Our principle location is at this place, Nauvoo (formerly Commerce[)] which is beautifully situated on the banks of the Mississippi, immediately above the lower rapids and is probably the best \& most beautiful site for a city on the River--- It has a gradual ascent from the river nearly a mile, then a fine level \& fertile Prairie, a situation in every respect adapted to commercial \& agricultural purposes; but like all other places on the river, is Sickly in summer.
					\hypertarget{EMS-1840-JS-8-AUG-1840-d}%
				The
					\marginnote{d}%
				number of inhabitants are nearly three thousand \& \sout{a} is fast increasing; if we are suffered to remain there is every prospect of its becoming one of the largest cities on the river if not in the western world, numbers have moved in from the Sea board and a few from the Islands of the Sea (Gr\supers{t.} Britain). It is our intention to commence the erection of some publick buildings next spring. We have purchased twenty thousand acres of land in the Iowa Territory oposite this place which is fast filling up with our people. \sout{It} I \sout{is my} desire \sout{that} all the Saints as well as all lovers of truth \& correct principles to come to this place as fast as possible as their circumstances will permit and endeavor by energy of action and a concentration of talent \&c \&c to effect those objects that are so dear to us.
					\hypertarget{EMS-1840-JS-8-AUG-1840-e}%
				Therefore
					\marginnote{e}%
				my general invitation is ``Let all that will, come'' and take of the poverty of Nauvoo freely. I should be disposed to give you a special invitation to come as early as possible believing you will be of great service to us, however you must make arrangements according to your circumstances \&c. Were it possible for you to come here this season to suffer affliction with the people of God no one will be more pleased or give you a more cordial welcome than myself

					\hypertarget{EMS-1840-JS-8-AUG-1840-f}%
				A
					\marginnote{f}%
				charter has been obtained from the Legislature for a Rail road from Warsaw being immediately below the rapids of the Mississippi to this place a distance of about Tewenty miles which if carried into opperation will be of incalculable advantage to this place as steam Boats can only asscend the rapids at a high stage of water. The soil is good and I should think not inferior to any in the state. Cropps are abundant in this section of the Country, and I think provisions will be reasonable. I should be very happy could I make arrangements to meet you in Springfield at the time you mention but cannot promise myself that pleasure; if I should not, probably you could make it convenient to come and pay us a visit here prior to your removal.

					\hypertarget{EMS-1840-JS-8-AUG-1840-g}%
				Elder
					\marginnote{g}%
				Rigdon is very sick, and has been for nearly twelve months with the fever and Ague which disease is very prevalent here at this time; at present he is not able to leave his room

				\begin{tightright}
					Yours \&c,

					Joseph Smith J\supers{r.}
				\end{tightright}

				J[ohn] C. Bennett M.D.

				P.S. Yours of the 30\supers{th.} is just received in which I am glad to learn of your increasing desire to unite yourself with a people ``that are every way spoken against'' and the anxiety you feel for our welfare for which you have my best feelings and I pray that my Heavenly Father will pour out his choicest blessings in this world and enable you by his grace to overcome the evils which are in the world that you may secure a blissful immortality in the world that is to come.

				\begin{tightright}
					J. S. J\supers{r.}
				\end{tightright}
				
		% -----
		\subsubsection{Minutes, 17 August 1840}
		% -----
			\label{EMS-1840-JS-17-AUG-1840}
			% \label{EMS-1840-JS-DAY-3LETTERMONTH-YEAR}
			
			\begin{description}
				\item[Print Sources]
			\end{description}

			\begin{tabular}{p{0.5in}p{1in}p{0.5in}p{1in}}
				JSP	 	& ---		& JSR 		& --- \\
				DC	 	& ---		& HC 		& --- \\
				HDDC 	& ---		& TCHC 		& --- \\
			\end{tabular}
			% HDDC = woodford's DC diss; JSR = Marquardt's JS Revelations; TCHC = Vogel HC
			
			\begin{description}
				\item[Online Access] \href{https://www.josephsmithpapers.org/paper-summary/minutes-17-august-1840/1}{Here}
				% \item[Analysis] \hyperref[XX]{Here}
			\end{description}
			
			\begin{description}
				\item[Background]
			\end{description}
			
				I need to include some of the background here.
			
				% It might also be useful to give a two sentence context of the period: e.g., "This speech was given in Nauvoo to a small congregation one month prior to the destruction of the Nauvoo Expositor. These and other tensions may have been on the prophet's mind when he gave the speech."
			
			\begin{description}
				\item[Original Source]
			\end{description}
			
				% brief description of original source material, with reference to JSP or somewhere else for more info
				% Background material: it might be helpful to have a note that says whether a given document was presented as a speech, was written in a journal, etc.
				% Also a comment here about how reliable the source might be, or what shape it is in, if there are large omissions or parts that are hard to understand, and so on.
			
			\begin{description}
				\item[Text]
			\end{description}
			
				...After the evidences were heard, Presidents John Smith, Milliam [William] Marks, Hyrum Smith spoke on the subject--- President Joseph Smith jr then spok[e] at some length showing the situation of the contending parties--- that there was in reality no cause of difference--- that they had better be reconciled without an action or vote of the Councils and hence forth live as brethren and never more mention their former difficulties \&c upon which John Patten replied that his hand was ready--- E. Fordham quickly responded to the same the parties then affectionately gave each other their hands in brotherly fellowship according to the council of President J. Smith j\supers{r} without an action of the Council...
				
		% -----
		\subsubsection{Discourse, 30 August 1840}
		% -----
			\label{EMS-1840-JS-30-AUG-1840}
			% \label{EMS-1840-JS-DAY-3LETTERMONTH-YEAR}
			
			\begin{description}
				\item[Print Sources]
			\end{description}

			\begin{tabular}{p{0.5in}p{1in}p{0.5in}p{1in}}
				JSP	 	& ---		& JSR 		& --- \\
				DC	 	& ---		& HC 		& --- \\
				HDDC 	& ---		& TCHC 		& --- \\
			\end{tabular}
			% HDDC = woodford's DC diss; JSR = Marquardt's JS Revelations; TCHC = Vogel HC
			
			\begin{description}
				\item[Online Access] \href{https://drive.google.com/file/d/1jxUt8QswSngeihPE6F9bUCEKcUZP2yYT/view?usp=share_link}{Here}
				% \item[Analysis] \hyperref[XX]{Here}
			\end{description}
			
			\begin{description}
				\item[Background]
			\end{description}
			
				% It might also be useful to give a two sentence context of the period: e.g., "This speech was given in Nauvoo to a small congregation one month prior to the destruction of the Nauvoo Expositor. These and other tensions may have been on the prophet's mind when he gave the speech."
			
			\begin{description}
				\item[Original Source]
			\end{description}
			
				It's from John Smith's journal, and I transcribed it myself on 12/10/22.
				
				The link above goes to a PDF file on my Drive, and the text is near the bottom of PDF 6.
			
				% brief description of original source material, with reference to JSP or somewhere else for more info
				% Background material: it might be helpful to have a note that says whether a given document was presented as a speech, was written in a journal, etc.
				% Also a comment here about how reliable the source might be, or what shape it is in, if there are large omissions or parts that are hard to understand, and so on.
			
			\begin{description}
				\item[Text]
			\end{description}
			
				30 sabbath Joseph Smith Jr continued his Discourse on Eternal Judgement and the Eternal Duration of matter
				
		% -----
		\subsubsection{Minutes, 5--6 September 1840}
		% -----
			\label{EMS-1840-JS-5-6-SEP-1840}
			% \label{EMS-1840-JS-DAY-3LETTERMONTH-YEAR}
			
			\begin{description}
				\item[Print Sources]
			\end{description}

			\begin{tabular}{p{0.5in}p{1in}p{0.5in}p{1in}}
				JSP	 	& ---		& JSR 		& --- \\
				DC	 	& ---		& HC 		& --- \\
				HDDC 	& ---		& TCHC 		& --- \\
			\end{tabular}
			% HDDC = woodford's DC diss; JSR = Marquardt's JS Revelations; TCHC = Vogel HC
			
			\begin{description}
				\item[Online Access] \href{https://www.josephsmithpapers.org/paper-summary/minutes-5-6-september-1840/1}{Here}
				% \item[Analysis] \hyperref[XX]{Here}
			\end{description}
			
			\begin{description}
				\item[Background]
			\end{description}
			
				I need to include some of the background here.
			
				% It might also be useful to give a two sentence context of the period: e.g., "This speech was given in Nauvoo to a small congregation one month prior to the destruction of the Nauvoo Expositor. These and other tensions may have been on the prophet's mind when he gave the speech."
			
			\begin{description}
				\item[Original Source]
			\end{description}
			
				% brief description of original source material, with reference to JSP or somewhere else for more info
				% Background material: it might be helpful to have a note that says whether a given document was presented as a speech, was written in a journal, etc.
				% Also a comment here about how reliable the source might be, or what shape it is in, if there are large omissions or parts that are hard to understand, and so on.
			
			\begin{description}
				\item[Text]
			\end{description}
			
					\hypertarget{EMS-1840-JS-5-6-SEP-1840-a}%
				September
					\marginnote{a}%
				5\supers{th} 1840 High Council met in council at the office of Joseph Smith j\supersu{r}

				Resolved. That Austin Cowles be appointed a member of this High Council, to fill the vacancy occasioned by the death of Elder Seymour Brunson.

				Joseph Smith j\supers{r} against Almon Bab[b]itt.

				The charge was predicated upon the authority of two letters, one from Thomas Burdock [Burdick], the other from Levi Richards and Oliver Granger accusing him (Babitt) with the following charges.

					\hypertarget{EMS-1840-JS-5-6-SEP-1840-b}%
				1\supers{st}
					\marginnote{b}%
				For stating that Joseph j\supersu{r} had extravagantly purchased three suits of cloths, while he was at Washington City and that Sidney Rigdon had purchased four suits at the same place, besides dresses and cloths for their families in profusion.

				2\supers{nd} For having stated that Joseph Smith j\supersu{r}--- Sidney Rigdon, and Elias Higbee had stated that they were worth \$\uline{100.000} each while they were at Washington \& that Joseph Smith j\supersu{r} had \sout{reported} <repeated> the same \sout{thing} statement while at Philadelphia and for stating that Oliver Granger had stated that he also, was worth as much as they (iE) \uline{100,000}

				3\supers{nd} For holding secret Council in the Lord's house, in Kirtland, and for locking the doors of the house, for the purpose of prohibiting certain brethren, in good standing, in the Church, from being in the Council. thereby depriving them the use of the house.

					\hypertarget{EMS-1840-JS-5-6-SEP-1840-c}%
				Two
					\marginnote{c}%
				were appointed to speak on the case, namely (7) Tho\supersu{s} Grover (8) A. Cowles.

				Council adjourned till the 6\supers{th} of Sep\supersu{t} at 2' o'clock

				Sept. 6\supers{th} 1840. Council met according to adjournment. When the evidences were all heard on the case pending and the council closed on both sides the parties spoke at length after which Pres\supersu{t} J. Smith jr withdrew the the charge and both parties were reconciled together things being adjusted to the satisfaction of both parties.

				Adjourned.

				\begin{tightright}
					Hosea Stout Clerk pro tem.
				\end{tightright}
				
		% -----
		\subsubsection{Letter to Saints Scattered Abroad, September 1840}
		% -----
			\label{EMS-1840-JS-SEP-1840}
			% \label{EMS-1840-JS-DAY-3LETTERMONTH-YEAR}
			
			\begin{description}
				\item[Print Sources]
			\end{description}

			\begin{tabular}{p{0.5in}p{1in}p{0.5in}p{1in}}
				JSP	 	& ---		& JSR 		& --- \\
				DC	 	& ---		& HC 		& --- \\
				HDDC 	& ---		& TCHC 		& --- \\
			\end{tabular}
			% HDDC = woodford's DC diss; JSR = Marquardt's JS Revelations; TCHC = Vogel HC
			
			\begin{description}
				\item[Online Access] \href{https://www.josephsmithpapers.org/paper-summary/letter-to-saints-scattered-abroad-september-1840/1}{Here}
				% \item[Analysis] \hyperref[XX]{Here}
			\end{description}
			
			\begin{description}
				\item[Background]
			\end{description}
			
				I need to include some of the background here.
			
				% It might also be useful to give a two sentence context of the period: e.g., "This speech was given in Nauvoo to a small congregation one month prior to the destruction of the Nauvoo Expositor. These and other tensions may have been on the prophet's mind when he gave the speech."
			
			\begin{description}
				\item[Original Source]
			\end{description}
			
				% brief description of original source material, with reference to JSP or somewhere else for more info
				% Background material: it might be helpful to have a note that says whether a given document was presented as a speech, was written in a journal, etc.
				% Also a comment here about how reliable the source might be, or what shape it is in, if there are large omissions or parts that are hard to understand, and so on.
			
			\begin{description}
				\item[Text]
			\end{description}
			
				\begin{tightcenter}
						\hypertarget{EMS-1840-JS-SEP-1840-a}%
					TO
						\marginnote{a}%
					THE SAINTS SCATTERED ABROAD.
				\end{tightcenter}

				\textsc{Beloved Brethren:}

				We address a few lines to the church of Jesus Christ, who have obeyed from the heart, that form of doctrine which has been delivered to them by the servants of the Lord, and who are desirous to go forward in the ways of truth and righteousness, and by obedience to the heavenly command, escape the things which are coming on the earth and secure to themselves an inheritance among the sanctified in the world to come.

				Having been placed in a very responsible situation in the church, we at all times feel interested in the welfare of the Saints and make mention of them continually in our prayers to our heavenly Father, that they may be kept from the evils which are in the world and ever be found walking in the path of truth.

					\hypertarget{EMS-1840-JS-SEP-1840-b}%
				The
					\marginnote{b}%
				work of the Lord in these last days, is one of vast magnitude and almost beyond the comprehension of mortals: its glories are past description and its grandeur insurpassable. It has been the theme which has animated the bosom of prophets and righteous men from the creation of this world down through every succeeding generation to the present time; and it is truly the dispensation of the fulness of times, when all things which are in Christ Jesus, whether in heaven or on the earth, shall be gathered together in him, and when all things shall be restored, as spoken of by all the holy prophets since the world began: for in it will take place the glorious fulfillment of the promises made to the fathers, while the displays of the power of the Most High will be great, glorious, and sublime.

					\hypertarget{EMS-1840-JS-SEP-1840-c}%
				The
					\marginnote{c}%
				purposes of our God are great, his love unfathomable, his wisdom infinite, and his power unlimited; therefore, the Saints have cause to rejoice and be glad, knowing that ``this God is our God forever and ever and he will be our guide unto death.''

				Having confidence in the power, wisdom and love of God, the Saints have been enabled to go forward through the most adverse circumstances, and frequently when to all human appearances nothing but death presented itself, and destruction, inevitable, has the power of God been manifest, his glory revealed, and deliverance effected; and the Saints, like the children of Israel who came out of the land of Egypt, and through the Red Sea, have sung an anthem of praise to his holy name: this has not only been the case in former ages, but in our own days, and within a few months, have we seen this fully verified.

					\hypertarget{EMS-1840-JS-SEP-1840-d}%
				Having,
					\marginnote{d}%
				through the kindness of our God, been delivered from destruction, and secured a location upon which we have again commenced opperations for the good of his people, we feel disposed to go forward and unite our energies for the upbuilding of the kingdom, and establishing the Priesthood in their fulness and glory.

				The work which has to be accomplished in the last days is one of vast importance, and will call into action the energy, skill, talent, and ability of the Saints, so that it may roll forth with that glory and majesty described by the prophets: and will consequently require the concentration of the Saints, to accomplish works of such magnitude and grandeur.

					\hypertarget{EMS-1840-JS-SEP-1840-e}%
				The
					\marginnote{e}%
				work of the gathering spoken of in the scriptures, will be necessary to bring about the glories of the last dispensation: It is probably unnecessary to press this subject on the Saints, as we believe the spirit of it is manifest, and its necessity obvious to every considerate mind; and every one zealous for the promotion of truth and righteousness, is equally so for the gathering of the Saints.

					\hypertarget{EMS-1840-JS-SEP-1840-f}%
				Dear
					\marginnote{f}%
				brethren feeling desirous to carry out the purposes of God, to which we have been called; and to be co-workers with him in this last dispensation: we feel the necessity of having the hearty co-operation of the Saints throughout this land, and upon the Islands of the sea; and it will be necessary for them to hearken to council, and turn their attention to the church, the establishment of the kingdom, and lay aside every selfish principle, every thing low, and groveling; and stand forward in the cause of truth, and assist to the utmost of their power, those to whom has been given the pattern and design; and like those who held up the hands of Moses, hold up the hands of those who are appointed to direct the affairs of the kingdom, so that they may be strengthened, and be enabled to prosecute their great designs and be instrumental in effecting the great work of the last days.

					\hypertarget{EMS-1840-JS-SEP-1840-g}%
				Believing
					\marginnote{g}%
				the time has now come when it is necessary to erect a house of prayer, a house of order, a house for the worship of our God; where the ordinances can be attended to agreably to his divine will, in this region of country; to accomplish which, considerable exertion must be made, means will be required; and as the work must be hastened in righteousness, it behooves the Saints, to weigh the importance of these things, in their minds, in all their bearings, and then take such steps as are necessary to carry them into operation; and arm themselves with courage, resolve to do all they can, and feel themselves as much interested, as though the whole labor depended on themselves alone; by so doing they will emulate the glorious deeds of the Fathers, and secure the blessing of heaven upon themselves and their posteri[t]y to the latest generation.

					\hypertarget{EMS-1840-JS-SEP-1840-h}%
				To
					\marginnote{h}%
				those who feel thus interested, and can assist in this great work, we say let them come to this place, by so doing they will not only assist in the rolling of the kingdom, but be in a situation where they can have the advantages of instruction from the presidency and other authorities of the church, and rise higher and higher in the scale of intel[li]gence, until they ``can comprehend with all Saints the length and breadth and debth and height, and know the love of God which passeth knowledge.''

					\hypertarget{EMS-1840-JS-SEP-1840-i}%
				Connected
					\marginnote{i}%
				with the building up of the kingdom, is the printing and circulation of the Book of Mormon, Doctrine and Covenants, Hymn book and the new translation of Scriptures, It is unnecessary to say any thing respecting these works; those who have read them, and who have drank of the stream of knowledge, which they convey, know how to appreciate them, and although fools may have them in derision, yet they are calculated to make men wise unto salvation, and sweep away the cobwebs of superstition of ages, throw a light on the proceedings of Jehovah which have already been accomplished and mark out the future in all its dreadful and glorious realities; those who have tasted the benefit derived from a study of those works, will undoubtedly vie with each other in their zeal for sending them abroad throughout the world, that every son of Adam may enjoy the same privileges and rejoice in the same truths.

					\hypertarget{EMS-1840-JS-SEP-1840-j}%
				Here
					\marginnote{j}%
				then, beloved brethren is a work to engage in worthy of arch-angels; a work which will cast into the shade the things which have heretofore been accomplished; a work which kings and prophets and righteous men, in former ages have sought, expected, and earnestly desired to see, but died without the sight: and well, will it be for those who shall aid in carrying into effect the mighty operations of Jehovah.

				\begin{tightright}
					By order of the first Presidency,

					R[obert] B. THOMPSON,

					Scribe.
				\end{tightright}

				\textit{Nauvoo, Sept.} 1840.
				
		% -----
		\subsubsection{Minutes and Discourse, 3--5 October 1840}
		% -----
			\label{EMS-1840-JS-3-5-OCT-1840}	
			% \label{EMS-1840-JS-DAY-3LETTERMONTH-YEAR}
			
			% --
			\paragraph{As Reported by Robert B. Thompson}
			% --
				\label{EMS-1840-JS-3-5-OCT-1840-RT}
				
				\begin{description}
					\item[Print Sources]
				\end{description}

				\begin{tabular}{p{0.5in}p{1in}p{0.5in}p{1in}}
					JSP	 	& ---		& JSR 		& --- \\
					DC	 	& ---		& HC 		& --- \\
					HDDC 	& ---		& TCHC 		& --- \\
				\end{tabular}
				% HDDC = woodford's DC diss; JSR = Marquardt's JS Revelations; TCHC = Vogel HC
				
				\begin{description}
					\item[Online Access] \href{https://www.josephsmithpapers.org/paper-summary/minutes-and-discourse-3-5-october-1840/1}{Here}
					% \item[Analysis] \hyperref[XX]{Here}
				\end{description}
				
				\begin{description}
					\item[Background]
				\end{description}
				
					I need to include some of the background here, and I need to check the JSP footnotes carefully. It is possible that there is other material that I can get and include.
				
					% It might also be useful to give a two sentence context of the period: e.g., "This speech was given in Nauvoo to a small congregation one month prior to the destruction of the Nauvoo Expositor. These and other tensions may have been on the prophet's mind when he gave the speech."
				
				\begin{description}
					\item[Original Source]
				\end{description}
				
					% brief description of original source material, with reference to JSP or somewhere else for more info
					% Background material: it might be helpful to have a note that says whether a given document was presented as a speech, was written in a journal, etc.
					% Also a comment here about how reliable the source might be, or what shape it is in, if there are large omissions or parts that are hard to understand, and so on.
				
				\begin{description}
					\item[Text]
				\end{description}
				
						\hypertarget{EMS-1840-JS-3-5-OCT-1840-RT-a}%
					...The
						\marginnote{a}%
					president arose and stated that there had been several depredations committed on the citizens of Nauvoo, and thought it expedient that a committee be appointed, to search out the offenders, and bring them to justice...
					
					The president then rose, and stated that it was necessary that something, should be done with regard to Kirtland, so that it might be built up; and gave it as his opinion, that the brethren from the east might gather there, and also, that it was necessary that some one should be appointed from this conference to preside over that stake...
					
					The president then spoke of the necessity of building a ``House of the Lord'' in this place.
					
						\hypertarget{EMS-1840-JS-3-5-OCT-1840-RT-b}%
					Whereupon
						\marginnote{b}%
					it was resolved, that the saints build a house for the worship of God, and that Reynolds Cahoon, Elias Higbee, and Alpheus Cutler, be appointed a committee to build the same.
					
					On motion. Resolved, that a commencement be made ten days from this date, and that every tenth day be appropriated for the building of said house...
					
					The clerk was then called upon to read the report of the presidency, in relation to the city plot. after which the president made some observations on the situation of the debts on the city plot and advised that a committee be appointed to raise funds to liquidate the same...
					
						\hypertarget{EMS-1840-JS-3-5-OCT-1840-RT-c}%
					President
						\marginnote{c}%
					Joseph Smith jr. then arose and delivered a discourse on the subject of baptism for the dead. which was listened to with considerable interest, by the vast multitude assembled...%
						\footnote{%
							% \label{}%
							See the accounts of this discourse below by \hyperref[EMS-1840-JS-3-5-OCT-1840-VK]{Vilate Murray Kimball} and \hyperref[EMS-1840-JS-3-5-OCT-1840-PW]{Phebe Carter Woodruff}.
							} 
					
					Elder R. B. Thompson after a few preliminary remarks, read an article on the priesthood, composed by president Joseph Smith jr, after which,
					
					Elder Babbitt delivered an excellent discourse on the same subject at some considerable length...
					
					The president then made some observations, and pronounced his benediction on the assembly...
				
			% --
			\paragraph{As Reported by Vilate Murray Kimball}
			% --
				\label{EMS-1840-JS-3-5-OCT-1840-VK}
				
				\begin{description}
					\item[Print Sources]
				\end{description}

				\begin{tabular}{p{0.5in}p{1in}p{0.5in}p{1in}}
					JSP	 	& ---		& JSR 		& --- \\
					DC	 	& ---		& HC 		& --- \\
					HDDC 	& ---		& TCHC 		& --- \\
				\end{tabular}
				% HDDC = woodford's DC diss; JSR = Marquardt's JS Revelations; TCHC = Vogel HC
				
				\begin{description}
					\item[Online Access] \phantom{.}
						\begin{compactdesc}
							\item[Original]	\href{https://drive.google.com/file/d/1jG62RBHi4fB_-S-IOJ6l-Cv036I9Pv9x/view?usp=sharing}{Here}
							\item[Transcription] \href{https://drive.google.com/file/d/1ZeE-U2ukUyB5DOPglEtZwTfwr_uEFq1B/view?usp=sharing}{Here}
						\end{compactdesc}
					% \item[Analysis] \hyperref[XX]{Here}
				\end{description}
				
				\begin{description}
					\item[Background]
				\end{description}
				
					% It might also be useful to give a two sentence context of the period: e.g., "This speech was given in Nauvoo to a small congregation one month prior to the destruction of the Nauvoo Expositor. These and other tensions may have been on the prophet's mind when he gave the speech."
				
				\begin{description}
					\item[Original Source]
				\end{description}
				
					PDF 3 of the original, PDF 53 of the transcription.
				
					% brief description of original source material, with reference to JSP or somewhere else for more info
					% Background material: it might be helpful to have a note that says whether a given document was presented as a speech, was written in a journal, etc.
					% Also a comment here about how reliable the source might be, or what shape it is in, if there are large omissions or parts that are hard to understand, and so on.
				
				\begin{description}
					\item[Text]
				\end{description}
				
						\hypertarget{EMS-1840-JS-3-5-OCT-1840-VK-a}%
					I
						\marginnote{a}%
					had writen to you a short time before, so I thought I would delay writeing again until after conference, which closed last Monday. We had the largest and most interesting conference that ever has been since the Church was organised. The people that attended it, were estimated at four thousand. Some thought there was more. Much business was done and many good instructions given. President Smith has opend a new and glorious subject of late which has caused quite a revival in the Church. That is being baptised for the dead. Paul speaks of it in first Corinthians 15th Chapter 29th vers. Josesph has received a more full explaination of it by Revelation.
						\hypertarget{EMS-1840-JS-3-5-OCT-1840-VK-b}%
					He
						\marginnote{b}%
					says it is the privilege of this church to be baptised for all their kinsfolks that have died before this Gospel came forth; even back to their great Granfather and Mother if they have ben personally acquainted with them. By so doing we act as agents for them, and give them the privilege of comeing forth in the first resurection. He says they will have the Gospel preached them them in Prison but there is no such things as spirrits being baptised. He does not wholey discard sisters Booths Vishon; says it was to show her the necessity of being Baptised. Since this order has ben preached here, the waters have ben continually troubled. During conference there were sometimes from eight to ten Elders in the river at a time baptiseing.
						\hypertarget{EMS-1840-JS-3-5-OCT-1840-VK-c}%
					There
						\marginnote{c}%
					is a perticlelar order that the Elders have to adminester in, and to preseurve this order it was President Smiths advise that it should not be attended to only in this place. I want to be baptised for my Mother. I calculated to wate until you come home but the last time Joseph spoke upon the subject he advised every one to be up and a doing and liberate their friends from bondage as quick as posable. So I think I shall go forward this week, as there is a number of the neighbors going forward.
						\hypertarget{EMS-1840-JS-3-5-OCT-1840-VK-d}%
					Some
						\marginnote{d}%
					have alredy ben baptised a number of times over. They have to be baptised and confirmed for one person before they can be baptised for another. Those that have no friends on the earth to be baptised for them can $\diamond\diamond$nd ministering spirits to whom so ever they will and make known their request. Thus you see there is a chance for all. Is not this a glorious doctrine? Surely the Gentiles will mock but we will rejoice in it
				
			% --
			\paragraph{As Reported by Phebe Carter Woodruff}
			% --
				\label{EMS-1840-JS-3-5-OCT-1840-PW}
				
				\begin{description}
					\item[Print Sources]
				\end{description}

				\begin{tabular}{p{0.5in}p{1in}p{0.5in}p{1in}}
					JSP	 	& ---		& JSR 		& --- \\
					DC	 	& ---		& HC 		& --- \\
					HDDC 	& ---		& TCHC 		& --- \\
				\end{tabular}
				% HDDC = woodford's DC diss; JSR = Marquardt's JS Revelations; TCHC = Vogel HC
				
				\begin{description}
					\item[Online Access] \phantom{.}
						\begin{compactdesc}
							\item[Original]	\href{https://drive.google.com/file/d/1HQHdsYg5vEQzLCv_edqMcre8PYZwVaSP/view?usp=sharing}{Here}
							\item[Transcription] \href{https://wilfordwoodruffpapers.org/documents/dbbbb430-5a1a-464d-97c3-fcfe2faf2612/page/b942a5ea-8f98-4b38-be54-2c2eaea3d957}{Here}
						\end{compactdesc}
					% \item[Analysis] \hyperref[XX]{Here}
				\end{description}
				
				\begin{description}
					\item[Background]
				\end{description}
				
					% It might also be useful to give a two sentence context of the period: e.g., "This speech was given in Nauvoo to a small congregation one month prior to the destruction of the Nauvoo Expositor. These and other tensions may have been on the prophet's mind when he gave the speech."
				
				\begin{description}
					\item[Original Source]
				\end{description}
				
					% brief description of original source material, with reference to JSP or somewhere else for more info
					% Background material: it might be helpful to have a note that says whether a given document was presented as a speech, was written in a journal, etc.
					% Also a comment here about how reliable the source might be, or what shape it is in, if there are large omissions or parts that are hard to understand, and so on.
				
				\begin{description}
					\item[Text]
				\end{description}
				
						\hypertarget{EMS-1840-JS-3-5-OCT-1840-PW-a}%
					brother
						\marginnote{a}%
					Joseph was expected to preach today on the Priest hood but his health would not admit of it so A. Babbit took the stand and delivered an interesting discourse upon the same subject.--- yesterday brother Joseph spoke a short time from 29 verse of the 15 chap. of first Cor. of which I will here after speak \& this afternoon brother L. White spoke on the same subject which was very interes ting...Now
						\hypertarget{EMS-1840-JS-3-5-OCT-1840-PW-b}%
					a
						\marginnote{b}%
					few words from Brother Joseph sermon on the living's being baptized for the dead that they may be judged according to men in the flesh; he has learned by revelation that those in this church may be baptized for any of their relatives who are dead and had not a \sout{privalege} privaledge of hearing \sout{it} this gospel even for their children, parents, bothers, sisters, grandparents, uncles, \& aunts,--- but not for acquaintances unless they send a mi nistering spirit to their friends on earth,--- as soon as they are baptized for their friends they are released from prison and they can claim them in the resurrection and bring them into the celestial kingdom---
						\hypertarget{EMS-1840-JS-3-5-OCT-1840-PW-c}%
					this
						\marginnote{c}%
					doctrin is cord<i>ally received by the church and they are going forward in multitudes, some are going to be baptized as many as 16 times--- they have to be baptized and confirmed separately for every friend--- once in one day--- How many can a spirit be baptized \sout{wh}--- why not deputize a friend on earth to do it for them--- John Wesly can receive this work but how can his spirit be baptize in water--- It is the privilege of the oldest one in a family to be baptized for their friends if they desire it but they can give it to another if they choose Brother Joseph makes this doctrin look verry plain and consisten---
						\hypertarget{EMS-1840-JS-3-5-OCT-1840-PW-d}%
					he
						\marginnote{d}%
					has been bringing strange things forward to the church this season--- strong meat--- he has delivered a course of lectures this season which were verry interesting--- I could not attend only apart of the time but often wished you could be present--- he says that the throne of God stands on the earth like this earth only it is cleansed \&c \&c--- says that this earth is the wickedest planate that ever God made and that its inhabitants \sout{are} will have the greatest glory if they overcome evil---
						\hypertarget{EMS-1840-JS-3-5-OCT-1840-PW-e}%
					and
						\marginnote{e}%
					that it was the largest when made but at different times there has been parts taken from it, for instance when the ten tribes were lost and says they will all be restored again.--- Paul says I die da<i>ly brother Joseph says that he actually died and was raised to life again--- says the beasts that he fought with at Epesus often killed him and he came to life again--- if so Willford why may not our little Sarah be raised to life again--- He, brother Joseph has brought forth many new things this season... 

						\hypertarget{EMS-1840-JS-3-5-OCT-1840-PW-f}%
					19th
						\marginnote{f}%
					Yesterday Sunday, we had an interesting sermon \sout{we had} \sout{an interesting} \sout{subjec} upon the living's being baptized for the dead it was made so plain that none could dispute it reasonably and I should be glad to be baptized for my grandparents and some others but want to wait for you to do it would it not be best, I had ra ther you would do it than any person living. 
				
		% -----
		\subsubsection{Report of the First Presidency, 4 October 1840}
		% -----
			\label{EMS-1840-JS-4-OCT-1840}
			% \label{EMS-1840-JS-DAY-3LETTERMONTH-YEAR}
			
			\begin{description}
				\item[Print Sources]
			\end{description}

			\begin{tabular}{p{0.5in}p{1in}p{0.5in}p{1in}}
				JSP	 	& ---		& JSR 		& --- \\
				DC	 	& ---		& HC 		& --- \\
				HDDC 	& ---		& TCHC 		& --- \\
			\end{tabular}
			% HDDC = woodford's DC diss; JSR = Marquardt's JS Revelations; TCHC = Vogel HC
			
			\begin{description}
				\item[Online Access] \href{https://www.josephsmithpapers.org/paper-summary/report-of-the-first-presidency-4-october-1840/1}{Here}
				% \item[Analysis] \hyperref[XX]{Here}
			\end{description}
			
			\begin{description}
				\item[Background]
			\end{description}
			
				I need to include some of the background here, and I need to check the JSP footnotes carefully. It is possible that there is other material that I can get and include.
			
				% It might also be useful to give a two sentence context of the period: e.g., "This speech was given in Nauvoo to a small congregation one month prior to the destruction of the Nauvoo Expositor. These and other tensions may have been on the prophet's mind when he gave the speech."
			
			\begin{description}
				\item[Original Source]
			\end{description}
			
				% brief description of original source material, with reference to JSP or somewhere else for more info
				% Background material: it might be helpful to have a note that says whether a given document was presented as a speech, was written in a journal, etc.
				% Also a comment here about how reliable the source might be, or what shape it is in, if there are large omissions or parts that are hard to understand, and so on.
			
			\begin{description}
				\item[Text]
			\end{description}
			
					\hypertarget{EMS-1840-JS-4-OCT-1840-a}%
				REPORT
					\marginnote{a}%
				FROM THE PRESIDENCY.

				The	Presidency of the church of Jesus Christ of Latter Day Saints, would respectfully report; that they feel rejoicing to meet the saints at another general conference and under circumstances as favorable as the present. Since our settlement in Illinois, we have for the most part been treated with courtesy and respect, and a feeling of kindness and of sympathy, has generally been manifested by all cfasses [classes] of the community, who with us, deprecate the conduct of those men, whose dark and blackning deeds, are stamped with everlasting infamy and disgrace.

					\hypertarget{EMS-1840-JS-4-OCT-1840-b}%
				The
					\marginnote{b}%
				contrast between our past and present situation is great. Two years ago, mobs were threatening, plundering, driving and murdering the saints. Our burning houses enlightened the canopy of heaven. Our women and children houseless and destitute, had to wander from place to place, to seek a shelter from the rage of persecuting foes.--- Now we enjoy peace, and can worship the God of heaven and earth without molestation. And expect to be able to go forward and accomplish the great and glorious work to which we have been called. Under these circumstances we feel to congratulate the saints of the Most High, on the happy and pleasing change in our circumstances, condition and prospects, and which those who shared in the perils and distresses, undoubtedly appreciate.
					\hypertarget{EMS-1840-JS-4-OCT-1840-c}%
				While
					\marginnote{c}%
				prayers and thanksgivings daily ascend to that God, who looked upon our distresses and delivered us from danger and death, and whose hand is over us for good. From the unpropitious nature of the weather, we hardly expected to behold so many of our friends on this occasion, in this however, we are agreeably disappointed, which gives us strong assurance that the saints are as zealous, untireing and energetic as ever in the great work of the last days; and gives us joy and consolation, and greatly encourages us, while contending with the difficulties which necessarily lie in our way.

					\hypertarget{EMS-1840-JS-4-OCT-1840-d}%
				Let
					\marginnote{d}%
				the brethren ever manifest such a spirit, and hold up our hands, and \textsc{we must, we will} go forward, the work of the Lord shall roll forth, the Temple of the Lord be reared, the Elders of Israel be encouraged; Zion be built up. And become the praise, the joy, and the glory of the whole earth; and the song of praise, glory, honor and majesty to him that setteth upon the throne, and to the Lamb forever and ever, shall reverberate from hill to hill, from mountain to mountain, from Island to Island and from continent to continent, and the kingdoms of this world become the kingdom of our God and his Christ.

					\hypertarget{EMS-1840-JS-4-OCT-1840-e}%
				We
					\marginnote{e}%
				are glad indeed to know that there is such a spirit of union existing throughout the churches, and at home and abroad; on this continent, as well as on the Islands of the sea, for by this principle and by a concentration of action shall we be able to carry into effect the purposes of our God.

				From the Elders abroad we receive the most cheering accounts; wherever the faithful laborer has gone forth reaping, sowing the seed of truth, he has returned with joy, bringing his sheaves with him; and the information we receive from all quarters is, that the laborers are few and that the harvest is great. Many wealthy and influential characters have embraced the gospel, so that not only will the poor rejoice in that they are exalted, but the rich in that they are made low.

					\hypertarget{EMS-1840-JS-4-OCT-1840-f}%
				The
					\marginnote{f}%
				calls to the southern states are indeed great, many places which a short time ago would think it a disgrace to give shelter to a Mormon, on account of the many false misrepresentations which were abroad, now desire to hear an Elder of the church of Latter Day Saints.

				On the Islands of the sea, viz. great Britain, there continues to be a steady flow of souls into the church--- branches have been organized in many large and populous cities and the whole land appears to be thirsting for the pure streams of knowledge and salvation. The Twelve have already printed a new edition of the Hymn book, and issue a monthly periodical in that land. Several families have already arrived here from England and a number more are on their way to this place, and are expected this fall.

					\hypertarget{EMS-1840-JS-4-OCT-1840-g}%
				If
					\marginnote{g}%
				the work roll forth with the same rapidity it has heretofore done, we may soon expect to see flocking to this place, people from every land and from every nation, the polished European, the degraded Hottentot, and the shivering Laplander. Persons of all languages, and of every tongue, and of every color; who shall with us worship the Lord of Hosts in his holy temple, and offer up their orisons in his sanctuary. It was in consideration of these things, and that a home might be provided for the saints, that induced us to purchase the present city for a place of gathering for the saints. and the extensive tract of land on the opposite side of the Mississippi. Although, the purchase at that time and under the peculiar conditions and circumstances of the church, appeared to many to be large and uncalled for; yet from what we now see, it is apparent to all, that we shall soon have to say. ``The place is too strait give us room that we may dwell.''

					\hypertarget{EMS-1840-JS-4-OCT-1840-h}%
				We
					\marginnote{h}%
				therefore hope that the brethren, who feel interested in the cause of truth, and desire to see the work of the gathering of Israel roll forth with power will aid us in liquidating the debts which are now owing, so that the inheritances may be secured to the church, and which eventually will be of great value. From the good spirit which is manifested on this occasion, the desire to do good, and the zeal for the honor of the church, inspires us with confidence that we shall not appeal in vain, but that funds will be forthcoming on this occasion, sufficient to meet the necessities of the case.

				It is with great pleasure that we have to inform the church that, another edition of the book of Mormon has been printed, and which is expected on from Cincinnatti, in a short time. And that arrangements are making for printing the book of Doctrine and Covenants Hymn book, etc. etc. So that the demand which may exist, for those works will soon be supplied.

					\hypertarget{EMS-1840-JS-4-OCT-1840-i}%
				In
					\marginnote{i}%
				conclusion we would say. Brethren and Sisters be faithful, be diligent, contend earnestly for the faith once delivered to the saints---let every man, woman and child realize the importance of the work, and act as if its success depended on their individual exertion alone, let them feel an interest in it, and then consider they live in a day, the contemplation of which animated the bosom of Kings, Prophets and Righteous men, thousands of years ago the prospect of which inspired their sweetest notes and most exalted lays' and caused them to break out in such rapturous strains as are recorded in the scriptures; and by and by, we shall have to exclaim in the language of Inspiration,

				``The Lord has brought again Zion---

				The Lord hath redeemed his people, Israel.''
			
		% -----
		\subsubsection{Instruction on Priesthood, circa 5 October 1840}
		% -----
			\label{EMS-1840-JS-CIR-5-OCT-1840}
			% \label{EMS-1840-JS-DAY-3LETTERMONTH-YEAR}
			
			\begin{description}
				\item[Print Sources]
			\end{description}

			\begin{tabular}{p{0.5in}p{1in}p{0.5in}p{1in}}
				JSP	 	& ---		& JSR 		& --- \\
				DC	 	& ---		& HC 		& --- \\
				HDDC 	& ---		& TCHC 		& --- \\
			\end{tabular}
			% HDDC = woodford's DC diss; JSR = Marquardt's JS Revelations; TCHC = Vogel HC
			
			\begin{description}
				\item[Online Access] \href{https://www.josephsmithpapers.org/paper-summary/instruction-on-priesthood-circa-5-october-1840/1}{Here}
				% \item[Analysis] \hyperref[XX]{Here}
			\end{description}
			
			\begin{description}
				\item[Background]
			\end{description}
			
				% It might also be useful to give a two sentence context of the period: e.g., "This speech was given in Nauvoo to a small congregation one month prior to the destruction of the Nauvoo Expositor. These and other tensions may have been on the prophet's mind when he gave the speech."
			
			\begin{description}
				\item[Original Source]
			\end{description}
			
				% brief description of original source material, with reference to JSP or somewhere else for more info
				% Background material: it might be helpful to have a note that says whether a given document was presented as a speech, was written in a journal, etc.
				% Also a comment here about how reliable the source might be, or what shape it is in, if there are large omissions or parts that are hard to understand, and so on.
			
			\begin{description}
				\item[Text]
			\end{description}
			
					\hypertarget{EMS-1840-JS-CIR-5-OCT-1840-a}%
				In
					\marginnote{a}%
				order to investigate the subject <of the Priesthood> so important to this as well as every succeeding generation, I shall proceed to trace the subject as far as I possibly can from the Old and new Testaments.

					\hypertarget{EMS-1840-JS-CIR-5-OCT-1840-b}%
				There
					\marginnote{b}%
				are two priesthoods spoken of in the scriptures, viz, the Melchisadeck and the Aaronic or Levitical Altho there are two priesthoods, yet the Melchisadeck priesthood comprehends the Aaronic or Levitical priesthood and is the Grand head and holds the hig[h]est Authority which pertains to the priesthood--- the keys of the Kingdom of God in all ages of the world to the latest posterity on the earth--- and is the channel through which all knowledge, doctrine, the plan of salvation and every important \sout{matter} truth is revealed from heaven.
					\hypertarget{EMS-1840-JS-CIR-5-OCT-1840-c}%
				Its
					\marginnote{c}%
				institution was prior to ``the foundation of this earth or the morning stars sang together or the sons of God shouted for joy,'' \sout{and} it is the highest and holiest priesthood and is after the order of the Son [of] God, and all other priesthoods are only parts, ramifications, powers and blessings belonging to the same and are held controlled and directed by it.
					\hypertarget{EMS-1840-JS-CIR-5-OCT-1840-d}%
				It
					\marginnote{d}%
				is the Channel through which the Almighty commenced revealing his glory at the beginning of the creation of this earth and through which he has continued to reveal himself to the children of men \sout{and} to the present time and through which he will make known his purposes to the end of time---

					\hypertarget{EMS-1840-JS-CIR-5-OCT-1840-e}%
				Commencing
					\marginnote{e}%
				with Adam who was the first man \sout{of} whom \sout{it} is spoken of in Daniel as being ``the Antient of days'' or in other words the first and oldest of all, the great grand progenitor of whom it is said in another place he is michael because he was the first and father of all, not only by progeny, but he was the first to hold the spiritual blessings, \sout{the plan} to whom was made known the plan of ordinances for the Salvation of his posterity unto the end, and to whom Christ was first revealed, and through whom christ has been revealed from heaven and will continue to be revealed from henceforth.
					\hypertarget{EMS-1840-JS-CIR-5-OCT-1840-f}%
				Adam
					\marginnote{f}%
				holds the keys of the dispensation of the fulness of times, i.e. the dispensation of all the times have been and will be revealed through him from the beginning to Christ and from Christ to the end of \sout{$\diamond\diamond$} all \sout{world} the dispensations that have to be reveald

					\hypertarget{EMS-1840-JS-CIR-5-OCT-1840-g}%
				Ephesians
					\marginnote{g}%
				1\supers{st.} chap 9 \& 10 verses. [``]Having made known unto us the mystery of his will, according to his good pleasure which he has purposed in himself that in the dispensation of the fulness of times he might gather together in one all things in Christ both which are in heaven and which are on earth in him'' Now the purpose in himself in the winding up scene of the last dispensation is, that all things pertaining to that dispensation should be conducted precisely in accordance with the preceeding dispensations,
					\hypertarget{EMS-1840-JS-CIR-5-OCT-1840-h}%
				And
					\marginnote{h}%
				again, God purposed in himself that there should not be an eternal fulness until every dispensation should be fulfilled and gathered together in one and that all things whatsoever that should be gathered together in one in those dispensations, unto the same fulness and eternal glory should be in Christ Jesus; therefore he set the ordinances to be the same for Ever and ever and set Adam to watch over them to reveal them from heaven to man or to send Angels to reveal them Heb 1 Chap. 16 verse. [``]Are they not all ministring spirits sent forth to minister to those who shall be heirs of salvation.'' These angels are under the direction of Michael or Adam who acts under the direction of Christ

					\hypertarget{EMS-1840-JS-CIR-5-OCT-1840-i}%
				From
					\marginnote{i}%
				the above quotation we learn that Paul perfectly understood the purpose of God in relation <to> his connexion with man, and that glorious and perfect order which he established in himself whereby he sent forth power revelations and glory. God will not acknowledge that which he has not called, ordained, and chosen.

					\hypertarget{EMS-1840-JS-CIR-5-OCT-1840-j}%
				In
					\marginnote{j}%
				the beginning God called Adam by his own voice See Genesis 3 Chap 9 \& 10 verses. [``]And the Lord called unto Adam and said unto him where art thou, and he said I heard thy voice in the garden and I was afraid because I was naked and hid myself.[''] Adam received commandments and instruction from God, this was the order from the begining: that he received revelations, Commandments, and ordinances at the begining is beyond the power of controversy, else, how did they begin to to offer sacrifices to God in an acceptable manner? [p. 3]

					\hypertarget{EMS-1840-JS-CIR-5-OCT-1840-k}%
				And
					\marginnote{k}%
				if they offered sacrifices they must be authorized by ordination. We read in Gen 4\supersu{th} chap. V.4 that [``]Abel brought of the firstlings of the flock and the fat thereof and the Lord had respect to Abel and to his offring.[''] And again Heb 11 chap 4 verse. [``]By Faith abel offered unto God a more excellent Sacrifice than Cain by which he obtained witness that he was righteous God testifying of his gifts and by it he being dead yet speaketh.['']
					\hypertarget{EMS-1840-JS-CIR-5-OCT-1840-l}%
				How
					\marginnote{l}%
				doth he yet speak? Why he magnified the priesthood which was confired [conferred] upon him and died a righteous man, and therefore has become a \sout{righteous} \sout{man} an angel of God by receiving his body from the dead, therefore holding still the keys of \sout{held} his dispensation and was sent down from heaven unto paul to minister consoleing words \& to comnin [communicate?] unto him a knowledge of the mysteries of Godliness
					\hypertarget{EMS-1840-JS-CIR-5-OCT-1840-m}%
				and
					\marginnote{m}%
				if this was not the case I would ask how did paul know so much about \sout{able} Abel and why should he talk about his speaking after he was dead. now that he spoke after he was dead must be, by being sent down out of heaven, to administer. This then is the nature of the priesthood, every man holding the presidency of his dispensation and one man holding the presidency of them all even Adam. and Adam receiving his Presidency and Authority from Christ, but cannot receive a fulness, untill Christ shall present, the Kingdom to the Father which shall be at the end of the last dispensation---

					\hypertarget{EMS-1840-JS-CIR-5-OCT-1840-n}%
				The
					\marginnote{n}%
				power, Glory, and blessings of the Priesthood could not continue with those who received ordination only as their righteousness continued, for Cain also being Authorized to offer sacrifice but not offering it in righteousness, therefore he was cursed. It signifies then, that the ordinances must be kept in the very way God has appointed, otherwise their priesthood will prove a cursing instead of a blessing. If Cain had fulfilled the law of righteousness as did Enoch he could have walked with God all the days of his life and never faild of a blessing.
					\hypertarget{EMS-1840-JS-CIR-5-OCT-1840-o}%
				Gen
					\marginnote{o}%
				\phantom{ }[\textit{blank}] [``]And Enoch walked with God after he begat Mathusalah 300 years and begat Sons and Daughters and all the days of Enoch were 365 years and Enoch walked with God and he was not for God took him.['']15 Now this Enoch God reserved unto himself that he should not die <at> \sout{and} that time and appointed unto him a ministry unto terrestrial bodies of whom there has been but little revealed, He is reserved also unto \sout{a} the presidency of a dispensation--- \sout{of} and more shall be said of him and terrestrial bodies in another treaties [treatise]

					\hypertarget{EMS-1840-JS-CIR-5-OCT-1840-p}%
				He
					\marginnote{p}%
				is a ministring Angel to minister \sout{for} to those who shall be heirs of salvation and appered unto Jude as Abel did unto paul. therefore Jude spoke of him 14 \& 15 verses in Jude. [``]and Enoch the Seventh revealed these sayings. Behold the Lord cometh with ten thousand of his saints[''] Paul was also aquainted with this Character and received instructions from him. Heb 11 Chap. 5 ver. [``]By Faith Enoch was translated that he should not see death, and was not found because God had translated him for before his translation he had this testimony that he pleased God. But without faith it is impossible to please God, for he that cometh to God must believe that he is, and that he is a revealer to those who diligently seek him,['']------

					\hypertarget{EMS-1840-JS-CIR-5-OCT-1840-q}%
				Now
					\marginnote{q}%
				the doctrine of translation is a power which belongs to this priesthood, there are many things which belong to the powers of the priesthood and the keys thereof that have been kept hid from before the foundation of the world. they are hid from the wise and prudent to be revealed in the last times. many may have supposed, that the--- doctrine of translation was a doctrine whereby men were taken immediately into the presence of God and into an Eternal fulness but this is a mistaken idea.
					\hypertarget{EMS-1840-JS-CIR-5-OCT-1840-r}%
				There
					\marginnote{r}%
				place of habitation is that of the terrestrial order and a place prepared for such characters, he held in reserve to be ministring Angels Unto many planets, and who as yet have not entered into so great a fulness as those who are resurrected from the dead, See Heb 11 Chap part of the 35 verse ``others were tortured not accepting deliverance that they might obtain a better resurrection[''] Now it was evident, that there was a better resurrection or else God would not have revealed it unto paul wherein then can it be said a better ressurrection?
					\hypertarget{EMS-1840-JS-CIR-5-OCT-1840-s}%
				This
					\marginnote{s}%
				distinction is made between the doctrine of the actual ressurrection and the doctrine of translation, the doctrine of translation obtains deliverance from the tortures and sufferings of the body but their existance will prolong as to their labors and toils of the ministry before they can enter in to so great a rest and glory, but on the other hand those who were tortured not accepting deliverance received an immediate rest from their labors,
					\hypertarget{EMS-1840-JS-CIR-5-OCT-1840-t}%
				See
					\marginnote{t}%
				Rev [\textit{blank}] [``]And I heard a voice from heaven saying blessed are the dead who die in the Lord for from henceforth they do rest from their labors and their works do follow them['']---23 They rest from their labors for a long time and yet their work is held in reserve for them, that they are permitted to do the same works after they receive a ressurection for their bodies, but we shall leave this subject and the subject of the terrestial bodies for another time in order to treat upon them more fully

					\hypertarget{EMS-1840-JS-CIR-5-OCT-1840-u}%
				The
					\marginnote{u}%
				next great grand patriarch who held the keys of the priesthood was Lamech See Gen 5 Chap 28 \& 29 verses--- [``]And Lamech lived 182 years and begat a Son and he called his name Noah saying this same shall comfort us concerning our work and the toil of our hands because of the ground which the Lord has Curst.[''] The priesthood continued from Lamech to Noah Gen 6 Chap 13 verse. [``]And God said unto Noah the end of all flesh is before me, for the earth is filled with violence through them, and behold I will destroy them with the earth,'' thus we behold the keys of this priesthood consisted in obtaining the voice of Jehovah that he talked with him in a familiare and friendly manner, that he continued to him the keys, the covenants, the power and the glory with which he blessed Adam at the beginning and the offring of sacrifice which also shall be continued at the last time.
					\hypertarget{EMS-1840-JS-CIR-5-OCT-1840-v}%
				for
					\marginnote{v}%
				all the ordinances and duties that ever have been required by the priesthood under the directions and commandments of the Almighty \sout{in the last dispensation at the end} \sout{thereof} in any of the dispensations, shall all be had in the last dispensations--- Therefore all things had under the Authority of the Priesthood At any former period shall be had again--- bringing to pass the restoration spoken of by the mouth of all the Holy prophets--- \sout{then} \sout{<Malach 3--- 3>} \sout{shall the Sons of Levi offer unto the Lord} \sout{an acceptable offering''} then shall the sons of Levi offer an acceptable sacrifice to the Lord Se[e] Malichi 3 Chap--- 3 \& 4 [``]And he shall sit as a refiners \sout{fire} and purifier of silver; and he shall purify the sons of Levi, and purge them as gold and silver, that they may offer unto the Lord['']

					\hypertarget{EMS-1840-JS-CIR-5-OCT-1840-w}%
				It
					\marginnote{w}%
				will be necessary here to make a few observations on the doctrine, set forth in the above quotation, As it is generally supposed that sacrifice was entirely done away, <when the great sacrif[ic]e was offered up---> and that there will be no necessity for the ordinance of sacrifice in future; but those who assert this, are certainly not aquainted with the duties, privileges and authority of the priesthood. <or with the prophets> The offering of sacrifice \sout{is} has ever been connected, and forms a part of the <duties of the> priesthood. It began \sout{which} with the priesthood and will be continued untill after the coming of christ from generation to generation------

					\hypertarget{EMS-1840-JS-CIR-5-OCT-1840-x}%
				We
					\marginnote{x}%
				frequently have mention made of the offering of Sacrifice by the servants of the most high in antient days prior to the law of moses, See [\textit{blank}] which ordinances will be continued when the priesthood is restored with all its Authority power and blessings. Elijah was the last Prophet that held the keys of this priesthood, and who will, before the last dispensation, restore the Authority and delive[r] the keys of this priesthood in order that all the ordinan[c]es may be attended to in righteousness. <It is true that the Savior had authority and power to bestow this blessing <but the sons of Levi were too predjudi[ced]>>

				And I will send Elijah the prophet before the great and terrible day of the Lord \&c \&c.

				Why send Elijah because he holds the keys of the Authority to administer in all the ordinances of the priesthood and without the Authority is given the ordinances could not be administered in righteousness.
				
					\hypertarget{EMS-1840-JS-CIR-5-OCT-1840-y}%
				It
					\marginnote{y}
				is a very prevalent opinion that \sout{in} the \sout{sacrifices} <sacrifices> \sout{of} <which> were offered were entirely consumed, this was not the case if you read Leviticus [\textit{blank}] Chap [\textit{blank}] Verses you will observe that the priests took a part as a memorial and offered it up before the Lord, while the remainder was kept for the \sout{benefit} <maintenance> of the priests--- So that the offerings and sacrifices are not all consumed upon the Alter, but the blood is sprinkled and the fat and certain other portions are consumed These sacrifices as well as every ordinance belonging to the priesthood will when the temple of the Lord shall be built <and the sons [of] Levi be purified> be fully restored and attended to \sout{then} \sout{the Sons of Levi shall be purified}. all their powers raniffications [ramifications], and blessings--- this ever \sout{was} did and will \sout{be} exist when the powers of the Melchisadc Priesthood are sufficiently manifest.
					\hypertarget{EMS-1840-JS-CIR-5-OCT-1840-z}%
				else
					\marginnote{z}%
				how can the restitution of all things spoken of by all the Holy prophets be brought to pass \sout{be brought to} \sout{pass}, It is not to be understood that, the law of Moses will be \sout{fully} established again with all its rights and <variety of \sout{ceremonies}> ceremonies <this has never been spoken off by the prophets> but those things which existed prior to Mose's day viz sacrifice, will be continued--- It may be asked by some what necesssity for Sacrifice since the great Sacrifice was offered? In answer to which if Repentance Baptism and faith \sout{were neccessary to Salvation} <existed> prior to the days of christ what necessity for them since that time------

				The priesthood has descended in a regular line from Father to Son through their succeeding generations

				See Book of Doctrine \& Covenants
				
		% -----
		\subsubsection{Letter to the Saints in Kirtland, Ohio, 19 October 1840}
		% -----
			\label{EMS-1840-JS-19-OCT-1840}
			% \label{EMS-1840-JS-DAY-3LETTERMONTH-YEAR}
			
			\begin{description}
				\item[Print Sources]
			\end{description}

			\begin{tabular}{p{0.5in}p{1in}p{0.5in}p{1in}}
				JSP	 	& ---		& JSR 		& --- \\
				DC	 	& ---		& HC 		& --- \\
				HDDC 	& ---		& TCHC 		& --- \\
			\end{tabular}
			% HDDC = woodford's DC diss; JSR = Marquardt's JS Revelations; TCHC = Vogel HC
			
			\begin{description}
				\item[Online Access] \href{https://www.josephsmithpapers.org/paper-summary/letter-to-the-saints-in-kirtland-ohio-19-october-1840/1}{Here}
				% \item[Analysis] \hyperref[XX]{Here}
			\end{description}
			
			\begin{description}
				\item[Background]
			\end{description}
			
				% It might also be useful to give a two sentence context of the period: e.g., "This speech was given in Nauvoo to a small congregation one month prior to the destruction of the Nauvoo Expositor. These and other tensions may have been on the prophet's mind when he gave the speech."
			
			\begin{description}
				\item[Original Source]
			\end{description}
			
				% brief description of original source material, with reference to JSP or somewhere else for more info
				% Background material: it might be helpful to have a note that says whether a given document was presented as a speech, was written in a journal, etc.
				% Also a comment here about how reliable the source might be, or what shape it is in, if there are large omissions or parts that are hard to understand, and so on.
			
			\begin{description}
				\item[Text]
			\end{description}
			
				\begin{tightright}
						\hypertarget{EMS-1840-JS-19-OCT-1840-a}%
					Nauvoo
						\marginnote{a}%
					Hancock County, Ills

					Oct 19\supers{th} 1840
				\end{tightright}

				To the Saints in

				Kirtland Ohio

				Dearly beloved brethren in the kingdom and patience of Jesus Christ.

				We take this opportunity of informing you that we yet remember the saints scattered about in the regions of Kirtland \& feel interested in their welfare as well as in that of the Saints at large.

				We have beheld with feelings peculiar to ourselves the situation of things in Kirtland and the numerous difficulties to which the Saints have been subjected by false friends as well as open enemies. All these circumstances have more or less engaged our attention from time to time.

					\hypertarget{EMS-1840-JS-19-OCT-1840-b}%
				We
					\marginnote{b}%
				likewise must complain of the \sout{stake} \sout{of} \sout{Kirtland} \sout{for} \sout{not} \sout{writing} \sout{to} \sout{us} \sout{the} brethren who are in office and authority in the Stake of Kirtland for not writing to us and making known their difficulties and their affairs from time to time so that they might be advised in matters of importance to the well being of said stake; but above all for not sending one word of consolation to us while we were in the hands of our enemies--- and thrust into dungeons.

				Some of our friends from various sections sent us letters which breathed a kind and sympathetick spirit, and which made our afflictions and sufferings [en]durable. All was silent as the grave no feelings of sorrow sympathy or affection to cheer the heart under the gloomy shades of affliction and trouble through which we had to pass.

					\hypertarget{EMS-1840-JS-19-OCT-1840-c}%
				Dear
					\marginnote{c}%
				bretheren could you realize that your bretheren were thus circumstancial and were to bear up under the weight of affliction and woe which was heaped upon them by their enemies and you stand unmoved and unconcerned!!! Where were the bowels of compasion, Where was the love which ought to characterize the Saints of the most high, did those high born and noble feelings lie dormant, or were you insensible of the treatment we received.

				However, we are disposed to leave these things to God and to futurity. and feel disposed to forget this coldness on the part of the Saints in Kirtland, and to look to the future with more pleasure than \sout{were} while we contemplate the past; and shall by the assistance of our Heavenly Father, take such steps as we think best calculated to promote the interest of the saints and for the promotion of truth \& righteousness, and the building up of the kingdom in these last days.

					\hypertarget{EMS-1840-JS-19-OCT-1840-d}%
				The
					\marginnote{d}%
				situation of Kirtland was brought before the general Conference holden at this place on the 3\supers{d.} instant, when it was resolved that Elder Almon Babbit[t] should be appointed to preside over the Stake of Kirtland and that he be privileged to choose his own counselors; We therefore hope that the Saints will hold up the hands of our beloved brother, and unite with him in endeavoring to promote the interest of the kingdom.

				It has been deemed prudent to advise the Eastern bretheren who desire to locate in Kirtland, to do so, consequently you may expect an increase of members in your stake, who probably will be but young in the faith and who will require kind treatment. We therefore hope the bretheren will feel interested in the welfare of the saints, and will use all their endeavors to promote the welfare of the bretheren who may think proper to take up their residence in that place.

					\hypertarget{EMS-1840-JS-19-OCT-1840-e}%
				If
					\marginnote{e}%
				you should put away from your midst all evil speaking, backbiting \& unge[ne]rous thoughts and feelings; humble yourselves and cultivate every principle of virtue \& Love then will the blessings of Jehovah rest upon you and you will yet see good and glorious days, peace will be within your gates, and prosperity in your borders which may our heavenly Father grant in the name of Jesus Christ is the prayer of your. in the bonds of the Covenant,

				\begin{tightright}
					Joseph Smith \supers{Jr.}

					Hyrum Smith 
				\end{tightright}

				P.S. Elder Almon Babbit is instructed to hold the keys of the \sout{Lords} House of the Lord. We therefore hope that they will be put into his hands as he will hold them for the benefit of the Church and subject to our instructions

				\begin{tightright}
					J Smith J\supers{r.}

					H. Smith
				\end{tightright}
				
		% -----
		\subsubsection{Bill to Incorporate the Church, 14 December 1840}
		% -----
			\label{EMS-1840-JS-14-DEC-1840}
			% \label{EMS-1840-JS-DAY-3LETTERMONTH-YEAR}
			
			\begin{description}
				\item[Print Sources]
			\end{description}

			\begin{tabular}{p{0.5in}p{1in}p{0.5in}p{1in}}
				JSP	 	& ---		& JSR 		& --- \\
				DC	 	& ---		& HC 		& --- \\
				HDDC 	& ---		& TCHC 		& --- \\
			\end{tabular}
			% HDDC = woodford's DC diss; JSR = Marquardt's JS Revelations; TCHC = Vogel HC
			
			\begin{description}
				\item[Online Access] \href{https://www.josephsmithpapers.org/paper-summary/bill-to-incorporate-the-church-14-december-1840/1}{Here}
				% \item[Analysis] \hyperref[XX]{Here}
			\end{description}
			
			\begin{description}
				\item[Background]
			\end{description}
			
				% It might also be useful to give a two sentence context of the period: e.g., "This speech was given in Nauvoo to a small congregation one month prior to the destruction of the Nauvoo Expositor. These and other tensions may have been on the prophet's mind when he gave the speech."
			
			\begin{description}
				\item[Original Source]
			\end{description}
			
				% brief description of original source material, with reference to JSP or somewhere else for more info
				% Background material: it might be helpful to have a note that says whether a given document was presented as a speech, was written in a journal, etc.
				% Also a comment here about how reliable the source might be, or what shape it is in, if there are large omissions or parts that are hard to understand, and so on.
			
			\begin{description}
				\item[Text]
			\end{description}
			
					\hypertarget{EMS-1840-JS-14-DEC-1840-a}%
				\phantom{ }
					\marginnote{a}%
				\sout{A} \sout{Bill} \sout{to} \sout{incorporate} \sout{the} \sout{Church} \sout{of} \sout{Jesus} \sout{Christ} \sout{of} \sout{Latter-Day} \sout{Saints}

				Sec. 1. Be it enacted by the People of the State of Illinois, represented in the General Assembly, That Joseph Smith, Hyrum Smith, Sidney Rigdon, William Marks, William Law, Elias Higbee, R[obert] B. Thompson, A[lanson] Ripley, Orson Pratt, D. C. [Don Carlos] Smith, Thomas Grover, D[imick B.] Huntington, A[rthur] Morrison, Charles [C.] Rich, T[heodore] Turley, H[enry] G. Sherwood, E[benezer] Robinson, G[eorge] W. Robinson, V[inson] Knight, and their associates and successors, the officers and members of the Church of Jesus Christ of Latter Day Saints, commonly called Mormons, are hereby created, constituted and declared to be, a body corporate and politic, with perpetual succession, with all and singular the powers, privilege\sout{s}, and prerogatives, of a corporation, by the name of the ``Church of Jesus Christ of Latter Day Saints,''
					\hypertarget{EMS-1840-JS-14-DEC-1840-b}%
				the
					\marginnote{b}%
				First Presidency of which, in conjunction with the General Assembly of said church, or the General Conference of said church, shall constitute the law-making department of the corporation for all secular purposes, with full power and authority to do all such acts as they may consider necessary for the welfare and prosperity of said church, and the said First Presidency, and their successors in office, shall be trustees in trust for the church, for the acquisition, regulation, or disposal, of all property, real, personal, or mixed, belonging to said church, or which may arise or accrue to the same; Provided, That no act shall be done repugnant to, or inconsistent with, the Constitution of the United States, or <the> \sout{of} \sout{this} \sout{state} <constitution and laws of this state>

					\hypertarget{EMS-1840-JS-14-DEC-1840-c}%
				Sec.
					\marginnote{c}%
				2. The corporators herein named shall meet, at the City of Nauvoo, on the 4th day of July next for the purpose of organization, and this act shall take effect, and be in force from, and after said day. [\textit{1/2 page blank}]

				[\textit{page 3 blank}]

				\begin{tightcenter}
					<to be Engrossed as amended

					M[eritt] L Covell Secty---

					strike out all after the Enacti[n]g

					clause and insert the following--->
				\end{tightcenter}
				
		% -----
		\subsubsection{Letter to Quorum of the Twelve, 15 December 1840}
		% -----
			\label{EMS-1840-JS-15-DEC-1840}
			% \label{EMS-1840-JS-DAY-3LETTERMONTH-YEAR}
			
			\begin{description}
				\item[Print Sources]
			\end{description}

			\begin{tabular}{p{0.5in}p{1in}p{0.5in}p{1in}}
				JSP	 	& ---		& JSR 		& --- \\
				DC	 	& ---		& HC 		& --- \\
				HDDC 	& ---		& TCHC 		& --- \\
			\end{tabular}
			% HDDC = woodford's DC diss; JSR = Marquardt's JS Revelations; TCHC = Vogel HC
			
			\begin{description}
				\item[Online Access] \href{https://www.josephsmithpapers.org/paper-summary/letter-to-quorum-of-the-twelve-15-december-1840/1}{Here}
				% \item[Analysis] \hyperref[XX]{Here}
			\end{description}
			
			\begin{description}
				\item[Background]
			\end{description}
			
				% It might also be useful to give a two sentence context of the period: e.g., "This speech was given in Nauvoo to a small congregation one month prior to the destruction of the Nauvoo Expositor. These and other tensions may have been on the prophet's mind when he gave the speech."
			
			\begin{description}
				\item[Original Source]
			\end{description}
			
				Check footnotes on the baptism for the dead stuff---there are other sources to look at.
			
				% brief description of original source material, with reference to JSP or somewhere else for more info
				% Background material: it might be helpful to have a note that says whether a given document was presented as a speech, was written in a journal, etc.
				% Also a comment here about how reliable the source might be, or what shape it is in, if there are large omissions or parts that are hard to understand, and so on.
			
			\begin{description}
				\item[Text]
			\end{description}
			
				\begin{tightright}
						\hypertarget{EMS-1840-JS-15-DEC-1840-a}%
					Nauvoo
						\marginnote{a}%
					Hancock Co, Ills. Dec\supersu{r}\supers{.} 15. 1840
				\end{tightright}

				Beloved Brethren.

				May Grace, Mercy, and Peace rest upon you, from God the Father and the Lord Jesus Christ.

				Having several communications laying before me, from my Brethren the ``Twelve'' some of which have ere this merited a a reply, but <from> the multiplicity of business which necessarily engages my attention I have delayed communicating to them, to the present time. Be assured, my beloved brethren, that I am no disinterested observer of the things which are transpiring on the face of the whole earth and amidst the general movements which are in progress, none is of more importance, than the glorious work in which you are now engaged,
					\hypertarget{EMS-1840-JS-15-DEC-1840-b}%
				and
					\marginnote{b}%
				consequently, I feel some anxiety on your account, that you may, by your virtue, faith, diligence, and charity, commend yourselves to one another, <to the Church of Christ> and <to> your Father which is in heaven, by whose grace you have been called to so holy a calling, and be enabled to perform the great and responsible duties which rest upon you. And I can assure you, that from the information I have received, I feel satisfied, that you have not been remiss, <in your duty> but that your diligence and faithfulness have been such, as must secure you the smiles of that God, whose servants you are, and the good will of the saints throughout the world.

					\hypertarget{EMS-1840-JS-15-DEC-1840-c}%
				The
					\marginnote{c}%
				spread of truth throughout England is certainly pleasing; the contemplation of which, cannot but afford feelings of no ordinary kind in the bosoms of those who have had to bear the heat and burthen of the day, and who were its firm supporters, and strenuous advocates, in infancy, while surrounded with circumstances the most unpropitious, and its destruction threatened on all hands. But like the gallant Bark, that has braved the storm unhurt, spreads her canvass to the breese, and nobly cuts her way through the yielding wave, more conscious than ever of the strength of her timbers and the experience and capabilities of her Captain, Pilate and crew.

					\hypertarget{EMS-1840-JS-15-DEC-1840-d}%
				It
					\marginnote{d}%
				is likewise very satisfactory to <my> mind, that there has been such a good understanding existing between you, and that the saints have <so> cheerfully, hearkened to council and vied with each other in their labors of love; and in the promotion of truth and righteousness; this is as it should be in the church of Jesus Christ. Unity is strength. ``How pleasant it is for brethren to dwell together in Unity \&c'' Let the saints of the most high, ever cultivate this principle, and the most glorious blessings must result, not only to them individually but to the whole church--- The order of the Kingdom will be maintained,--- Its officers respected, and its requirements readily and cheerfully obeyed.
					\hypertarget{EMS-1840-JS-15-DEC-1840-e}%
				Love
					\marginnote{e}%
				is one of the leading characteristics of Deity, and ou[gh]t to be manifested by those who aspire to be the Sons of God. A man filled with the love of God, is not content with blessing his family alone, but ranges through the world, anxious to bless the whole of the human family. This has been your feelings and caused you to forego the pleasures of home, that you might be a blessing to others, who are candidates for immortality \sout{but} and who were strangers to the principals of truth and for so doing I pray that Heaven's choicest blessings may rest upon you.

					\hypertarget{EMS-1840-JS-15-DEC-1840-f}%
				Being
					\marginnote{f}%
				requested to give my advice respecting the propriety of your returning in the spring, I will do so willingly. I have reflected on the subject some time and am of the opinion that it would be wisdom in you to make preparations to leave the scene of your labors in the spring. Having carried the testimony to that land, and numbers having received it, consequently the leaven can now spread, without your being obliged to stay. Another thing, there has been some whisperings of the spirit; that there will be some agitation, some excitement, and some trouble in the land in which you are now laboring. I would therefore say in the mean time be diligent, organize the churches and let every one stand in his proper place, so that those who cannot come with you in the spring may not be left as sheep without shepherds.

					\hypertarget{EMS-1840-JS-15-DEC-1840-g}%
				I
					\marginnote{g}%
				would likewise observe that inasmuch as this place has been appointed for the gathering of the saints, it is necessary that it should be attended to, in the order which the Lord intended it should; to this end I would say that as there are great numbers of the saints in England, who are extremely poor and not accustomed to the farming business, who must have certain preparations made for them before they can support themselves in this country, therefore to prevent confusion and disappointment when they arrive here, let those men who are accustomed to making machinery and those who can command a capital even if it be but small, come here as soon as convenient and put up machinery and make such other preparations as may be necessary, so, that when the poor come on they may have employment to come to. This place has advantages for \sout{a} manufacturing and commercial purposes which but very few can boast of; and by establishing Cotton Factories, Founderies, Potteries \&c \&c would be the means of bringing in wealth and raising it to a very important elevation. I need not occupy more space on this subject as its reasonableness must be obvious to every mind.
					\hypertarget{EMS-1840-JS-15-DEC-1840-h}%
				In
					\marginnote{h}%
				my former epistle I told you my mind respecting the printing of the Book of Mormon. Hymn Book \&c \&c--- I have been favored by receiving a Hymn Book from you and as far as I have examined it I highly approve of it and think it to be a very valuable collection. I am informed that the Book of Mormon is likewise printed, which I am glad to hear, and should be pleased to hear that it was printed in all the different Languages of the earth. You can use your own pleasure respecting printing the Book of Doctrine \& Covenants, if there is a great demand for them, I have not any objections, but would rather encourage it.

					\hypertarget{EMS-1840-JS-15-DEC-1840-i}%
				I
					\marginnote{i}%
				am happy to say, that as far as I have been made acquainted with your movements, I have been perfectly satisfied that they have been in wisdom, and I have no doubt but the spirit of the Lord has directed you. and this proves to my mind that you have been humble, and your desires have been, for the salvation of your fellow man, and not your own aggrandizement and selfish interest. As long as the saints manifest such a disposition their councils will be approved of, and their exertions crowned with success. There are many things of minor importance, on which you ask council, but which I think you will be perfectly able to decide upon as \sout{I} you are more conversant with the peculiar circumstances than I \sout{can} am, and I feel great confidence in your united wisdom, therefore you will excuse me for not entering into detail. If I should see any thing that was wrong, I should take the priviledge of making known my mind to you and pointing out the evil.

					\hypertarget{EMS-1840-JS-15-DEC-1840-j}%
				If
					\marginnote{j}%
				Elder Parley P Pratt should wish to remain in England, for some time longer than the rest of the Twelve, he will feel himself at liberty to do so; as he his family are with him Consequently his circumstances are somewhat different to the rest, and likewise it is necessary that some one should remain who is conversant with the rules, regulations \&c \& of the church And continue the paper which is published; consequently taking all these things into consideration I would not press upon Brother Pratt to return in the spring.

				I am happy to inform you that we are prospering in this place, and that the saints are more healthy than formerly, and from the decrease of sickness this season, when compared with the last, I am led to the conclusion that this, must eventually become a healthy place.

					\hypertarget{EMS-1840-JS-15-DEC-1840-k}%
				There
					\marginnote{k}%
				are at present about 3000 inhabitants in Nauvoo, and numbers are flocking in daily; severeal stakes have been set off in different parts of the \sout{country} county, which are in prospering circumstances. Provisions are much lower than when you left. Flour is worth about four dollars per barrel, corn \sout{25} 20 cents per bushel; Pottatoes about 20 cents. and other things in about the same proportion. There has been a very plentiful harvest indeed, throughout the Union.

				You will observe by the ``Times \& Season'' that we are about building a Temple for the worship of our God in this place; preparations are now making, every tenth day is devoted by the brethren here, for quarrying rock \&c \&. we have secured one of the most lovely sites for it that there is in this region of Country. It is expected to be considerably larger and on <a> more magnificent scale than the one in Kirtland and which will undoubtedly attract the attentio[n] of the great men of the <earth> We have a bill before the Legislature for the incorporation of the City of Nauvoo for the establishment of a Seminary and other purposes, which I expect will pass in a short time.

					\hypertarget{EMS-1840-JS-15-DEC-1840-l}%
				You
					\marginnote{l}%
				will also have received intelligence of the death of my Father, which event altho painful to the family and to the church generally, yet the sealing testimony of the truth of the work of the Lord was indeed satisfactory; the particulars of his death \&c you will find in the Sep\supersu{r}\supers{.} number of the ``Times and seasons'' Brother Hyrum [Smith] succeeds him as patriarch of the Church, according to his last directions and benedictions.

				Several persons of emminece [eminence] and distinction in society, have joined the Church, and become obedient to the faith, and I am happy to inform you that the work is spreading very fast on this continent, some of the Brethren are now in New Orleans, and we expect to have a gathering from the South.

					\hypertarget{EMS-1840-JS-15-DEC-1840-m}%
				I
					\marginnote{m}%
				have had the pleasure of wel[c]oming about One hundred of the Brethren from England, who came with Elder [Theodore] Turley, the remainder I am informed stop[p]ed in Kirtland, not having means to get any further. I think \sout{they can} those that came here did not take the best possible rout, or the least expensive. Most of the brethren have obtained employment of one kind or another and appear tolerably well contented and seem disposed to hearken to council. Brothers Robinson \& Smith lately had a letter from Elders [Heber C.] Kimball, [George A.] Smith \& [Wilford] Woodruff in London which gave us information of the commencement of the work of the Lord in that City, which I was glad to hear. I am likewise informed the Elders have gone to Austrailia \& to the East Indies I feel deserieous that every providential opening of that kind should be filled, and that you should prior to your leaving England, send the gospel into as many parts as you possibly can.

					\hypertarget{EMS-1840-JS-15-DEC-1840-n}%
				Beloved
					\marginnote{n}%
				brethren, you must be aware in some measure of my feelings when I contemplate the great work which is now rolling on, and the relationship which I sustain to it; while it is extending to distant lands, and islands, and thousands are embracing it, I realize in some measure my responsibility and the need I have of support from above, and wisdom from on high; that I may be able to teach this people, which have now become a great people, the principles of righteousness, and lead them agreeably to the will of heaven so that they may be perfected and prepared to meet the Lord Jesus Christ, when he shall appear. in great glory. Can I rely on your prayers to your \uline{heavenly} Father in my behalf? and on the prayers of all my brethren \& sisters in England? (whom having not seen yet I love) that I may be enabled to escape every stra[ta]gem of satan, surmount every difficulty, and bring this people, to the enjoyment of those blessings, which are reserved for the righteous I ask this at your and their hands in the name of Jesus Christ.

					\hypertarget{EMS-1840-JS-15-DEC-1840-o}%
				Let
					\marginnote{o}%
				the saints remember that great things depend on their individual exertion, and that they are called to be co-workers with us and the holy spirit in accomplishing the great works of the last days, and in consideration of the extent, the blessings, and the glories of the same let every selfish feeling be not only buried, but anihalated, and let love to God and man, predominate and reign triumphant in every mind, that their hearts may become like unto Enoch's of old so that they may comprehend all things, present, past, and future, and ``come behind in no gift waiting for the coming of the Lord Jesus Christ''.
					\hypertarget{EMS-1840-JS-15-DEC-1840-p}%
				The
					\marginnote{p}%
				work in which we are unitedly engaged in, is one of no ordinary kind, the enemies we have to contend against are subtle and well skilled in manuvering, it behoves us then to be on the alert, to concentrate our energies, and that the best feelings should exists in our midst, and then by the help of the Almighty we shall go on from victory to victory and from conquest unto conquest, our evil passions will be subdued, our predjudices depart, we shall find no room in our bosoms for hatred, vice will hide its deformed head, and we shall stand approved in the sight of heaven and be acknowledged ``the Sons of God'' Let us realize that we are not to live to ourselves but to God by so doing the greatest blessings will rest upon us both in time and in Eternity.

					\hypertarget{EMS-1840-JS-15-DEC-1840-q}%
				I
					\marginnote{q}%
				presume the doctrine of ``Baptizm for the dead'' has ere this reached your ears, and may have raised some inquiries in your mind respecting the same. I cannot in this letter give you all the information you may desire on the subject, but aside from my knowledg[e] independant of the Bible, I would say, that this was certainly practized by the antient churches And. St Paul endeavors to prove the doctrine of the ressurrection from the same, and says ``else what shall they do who are baptized for the dead[''] \&c \&c. I first mentioned the doctrine in public while preaching the funeral sermon of Bro [Seymour] Brunson, and have since then given general instructions to the Church on the subject.
					\hypertarget{EMS-1840-JS-15-DEC-1840-r}%
				The
					\marginnote{r}%
				saints have the priviledge of being baptized for those of their relatives who are dead, who they feel to believe would have embraced the gospel if they had been priviledged with hearing it, and who have received the gospel in the spirit through the instrumentality of those who may have been commissioned to preach to them while in the prison. Without enlarging on the subject you will undoubtedly see its consistancy, and reasonableness, and presents the the gospel of Christ in probably a more enlarged scale than some have received it. But as the performance of this right is more particularly confined to this place it will not be necessary to enter into particulars, at the same time I allways feel glad to give all the information in my power, but my space will not allow me to do it.

					\hypertarget{EMS-1840-JS-15-DEC-1840-s}%
				We
					\marginnote{s}%
				had a letter from Elder [Orson] Hyde a few days ago, who is in New Jersey, and is expecting to leave for England as soon as Elder [John E.] Page reaches him. He requested to know in his letter if converted Jews are to go to Jerusalem or to come to Zion. I therefore wish you to inform him that converted Jews must come here. If Elder Hydes \& Pages testimony to the Jews at Jerusalem should be received then they may know ``that the set time hath come'': I will write more particular instructions to them afterwards. Your

				Your families are well and generally in good spirits, and bear their privations with christian fortitude and patience.

				Brother [Willard] Richards' question respecting arriving in the spring is answered I shall be very happy to see him \& his family \& likewise Brother [Joseph] Fielding, tell him that Bro [Robert B.] Thompson is making preparations for his coming.

					\hypertarget{EMS-1840-JS-15-DEC-1840-t}%
				With
					\marginnote{t}%
				respect to the rout best to be taken I think you will be better able to <give> advise than myself. But I would not advise coming round by the lakes. And it would not be prudent to come via New Orleans in the sickly season. but in the spring or fall or winter it might do. Give my kind love to all the brethren, and sisters, and tell them I should have been pleased to have come over to England to see them, but am afraid that I shall be under the necessity of remaining here for some time, therefore I give them a pressing invitation to come and see me.

				I am D\supers{r} Brethren Yours, Affectionately \textbf{Joseph Smith}

				To the Travelling High Council and Elders of the Church of Jesus Christ of L.D.S. in Great Britain

				Do not understand me to say that all the Elders are to come with you, as it will be necessary for some to stay. J.

				To the ``Twelve''
				
		% -----
		\subsubsection{Discourse, December 1840, as Reported by William Clayton}
		% -----
			\label{EMS-1840-JS-DEC-1840}
			% \label{EMS-1840-JS-DAY-3LETTERMONTH-YEAR}
			
			\begin{description}
				\item[Print Sources]
			\end{description}

			\begin{tabular}{p{0.5in}p{1in}p{0.5in}p{1in}}
				JSP	 	& ---		& JSR 		& --- \\
				DC	 	& ---		& HC 		& --- \\
				HDDC 	& ---		& TCHC 		& --- \\
			\end{tabular}
			% HDDC = woodford's DC diss; JSR = Marquardt's JS Revelations; TCHC = Vogel HC
			
			\begin{description}
				\item[Online Access] \href{https://www.josephsmithpapers.org/paper-summary/discourse-december-1840-as-reported-by-william-clayton/1}{Here}
				% \item[Analysis] \hyperref[XX]{Here}
			\end{description}
			
			\begin{description}
				\item[Background]
			\end{description}
			
				% It might also be useful to give a two sentence context of the period: e.g., "This speech was given in Nauvoo to a small congregation one month prior to the destruction of the Nauvoo Expositor. These and other tensions may have been on the prophet's mind when he gave the speech."
			
			\begin{description}
				\item[Original Source]
			\end{description}
			
				% brief description of original source material, with reference to JSP or somewhere else for more info
				% Background material: it might be helpful to have a note that says whether a given document was presented as a speech, was written in a journal, etc.
				% Also a comment here about how reliable the source might be, or what shape it is in, if there are large omissions or parts that are hard to understand, and so on.
			
			\begin{description}
				\item[Text]
			\end{description}
			
				\begin{tightcenter}
					A key by Joseph Smith Dec\supersu{m} 1840---W. C. [William Clayton]
				\end{tightcenter}
				
				If an Angel or spirit appears offer him your hand; if he is a spirit from God he will stand still and not offer you his hand. If from the Devil he will either shrink back from you or offer his hand, which if he does you will feel nothing, but be deceived.
				
				\uline{A} \uline{good} \uline{Spirit} \uline{will} \uline{not} \uline{deceive.}
				
				Angels are beings who have bodies and appear to men in the form of man.